%-%-%-%-%-%-%-%-%-%-%-%-%-%-%-%-%-%-%-%-%-%-%-%-%-%-%-%-%-%-%-%-%-%-%-%-%-%-%
%%% MAIN DOCUMENT %%%

%-%-%-%-%-%-%-%-%-%-%-%-%-%-%-%-%-%-%-%-%-%-%-%-%-%-%-%-%-%-%-%-%-%-%-%-%-%-%

\documentclass[12pt, oneside]{book}


%%% AESTHETICS %%%
%-%-%-%-%-%-%-%-%-%-%-%-%-%-%-%-%-%-%-%-%-%-%-%-%-%-%-%-%-%-%-%-%-%-%-%-%-%-%


%%% Dimensions and Spacing %%%
\usepackage[margin=1in]{geometry}
\usepackage{setspace}
\linespread{1}


%%% Define new colors %%%
\usepackage{xcolor}
\definecolor{orangehdx}{rgb}{0.96, 0.51, 0.16}

% Normal colors
\definecolor{xred}{HTML}{BD4242}
\definecolor{xblue}{HTML}{4268BD}
\definecolor{xgreen}{HTML}{52B256}
\definecolor{xpurple}{HTML}{7F52B2}
\definecolor{xorange}{HTML}{FD9337}
\definecolor{xdotted}{HTML}{999999}
\definecolor{xgray}{HTML}{777777}
\definecolor{xcyan}{HTML}{80F5DC}
\definecolor{xpink}{HTML}{F690EA}
\definecolor{xgrayblue}{HTML}{49B095}
\definecolor{xgraycyan}{HTML}{5AA1B9}

% Dark colors
\colorlet{xdarkred}{red!85!black}
\colorlet{xdarkblue}{xblue!85!black}
\colorlet{xdarkgreen}{xgreen!85!black}
\colorlet{xdarkpurple}{xpurple!85!black}
\colorlet{xdarkorange}{xorange!85!black}
\definecolor{xdarkcyan}{HTML}{008B8B}
\colorlet{xdarkgray}{xgray!85!black}

% Very dark colors
\colorlet{xverydarkblue}{xblue!50!black}

% Document-specific colors
\colorlet{normaltextcolor}{black}
\colorlet{figtextcolor}{xblue}

% Enumerated colors
\colorlet{xcol0}{black}
\colorlet{xcol1}{xred}
\colorlet{xcol2}{xblue}
\colorlet{xcol3}{xgreen}
\colorlet{xcol4}{xpurple}
\colorlet{xcol5}{xorange}
\colorlet{xcol6}{xcyan}
\colorlet{xcol7}{xpink!75!black}

% Blue-Purple (should just used colorbrewer...)
\definecolor{xrainbow0}{HTML}{e41a1c}
\definecolor{xrainbow1}{HTML}{a24057}
\definecolor{xrainbow2}{HTML}{606692}
\definecolor{xrainbow3}{HTML}{3a85a8}
\definecolor{xrainbow4}{HTML}{42977e}
\definecolor{xrainbow5}{HTML}{4aaa54}
\definecolor{xrainbow6}{HTML}{629363}
\definecolor{xrainbow7}{HTML}{7e6e85}
\definecolor{xrainbow8}{HTML}{9c509b}
\definecolor{xrainbow9}{HTML}{c4625d}
\definecolor{xrainbow10}{HTML}{eb751f}
\definecolor{xrainbow11}{HTML}{ff9709}

%------- %
% XHFILL %
%------- %


%%% Chapter Headings %%%
\usepackage[Glenn]{fncychap}
\ChTitleVar{\bfseries\scshape\color{black}}
\ChNumVar{\large\selectfont\color{black}}
\ChNameVar{\large\color{black}}
\usepackage{xpatch}
\xpatchcmd\DOCH
{\mghrulefill}{\color{orangehdx}\mghrulefill}
{}{\PatchFailed}
\xpatchcmd\DOTI
{\mghrulefill}{\color{orangehdx}\mghrulefill}
{}{\PatchFailed}
\xpatchcmd\DOTIS
{\mghrulefill}{\color{orangehdx}\mghrulefill}
{}{\PatchFailed}

%% Change Chapter Heading Placement %%
\usepackage{etoolbox}
\makeatletter
\patchcmd{\@makechapterhead}{\vspace*{50\p@}}{\vspace*{-20\p@}}{}{}
\patchcmd{\@makeschapterhead}{\vspace*{50\p@}}{\vspace*{-20\p@}}{}{}
\patchcmd{\DOTI}{\vskip 80\p@}{\vskip 40\p@}{}{}
\patchcmd{\DOTIS}{\vskip 40\p@}{\vskip 0\p@}{}{}
\makeatother

\renewcommand{\thesection}{\thechapter.\arabic{section}} %% Chapter.Section Numbering


%%% FIGURES %%%
\usepackage{graphicx}  
\graphicspath{ {images/} }  
% \numberwithin{figure}{section}
\usepackage{float}
\usepackage{caption}

%%% Hyperlinks %%%
\usepackage{hyperref}
\definecolor{horange}{HTML}{f58026}
\hypersetup{
	colorlinks=true,
	linkcolor=horange,
	filecolor=horange,      
	urlcolor=horange,
}


%% Headers and Footers %%
\usepackage{fancyhdr} % This should be set AFTER setting up the page geometry
\pagestyle{fancy} % options: empty , plain , fancy
\renewcommand{\chaptermark}[1]{\markboth{\thechapter.\ #1}{}}
\fancyhead[R]{\leftmark}
\fancyhead[L]{Hendrix College}
\fancyhead[C]{\includegraphics[height=.50cm]{images/small logo.png}}
\usepackage{xpatch}
\xpretocmd\headrule{\color{orangehdx}}{}{\PatchFailed}
\setlength{\footskip}{0.5in}
\setlength{\headheight}{15pt} % Adjust accordingly if you have a tall header


%%%%% Colored Boxes %%%%%
\usepackage{tcolorbox}
\newcounter{BoxCounter}


\newtcolorbox{note}{colback=black!15!white, colframe=black, boxrule=0.5mm, arc=4mm, left=4mm, right=4mm, top=2mm, bottom=2mm}

%% Theorems %%
\newtcolorbox[use counter=BoxCounter, number within=section]{theorem}[1][]{%
    rounded corners,
    colframe=xblue, %color of frame
    colback=xblue!15, %color of background
    coltitle=white, %color of title text
    label={thm:\theBoxCounter},
    title={\large \textbf{Theorem \theBoxCounter}},
    #1
}

\newtcolorbox[use counter=BoxCounter, number within=section]{ntheorem}[2][]{%
    rounded corners,
    colframe=xblue, %color of frame
    colback=xblue!15, %color of background
    coltitle=white, %color of title text
    label={thm:#2},
    title={\large \textbf{Theorem \theBoxCounter: #2}},
    #1
}


\newcommand{\numrefthm}[1]{\hyperref[thm:#1]{Theorem #1}}

\newcommand{\namrefthm}[1]{\hyperref[thm:#1]{#1}}

% %% This command needs to exist for the square root of 2 theorem to work

% \newtcolorbox[]{sqrtTh}[2][]{%
%     rounded corners,
%     colframe=xblue, %color of frame
%     colback=xblue!15, %color of background
%     coltitle=white, %color of title text
%     label={thm:#2},
%     title={\large \textbf{Theorem: \texorpdfstring{The Irrationality of $\sqrt{2}$}{The Irrationality of sqrt(2)}}},
%     #1
% }


% Create the custom definition environment
\newtcolorbox[use counter=BoxCounter, number within=section]{definition}[1][]{%
    rounded corners,
    colframe=xred,
    colback=xred!15,
    coltitle=white,
    label={def:\theBoxCounter},
    title={\large \textbf{Definition \theBoxCounter}},
    #1
}

% Command to link specific words within the definition
\newcommand{\defword}[2]{\hyperref[def:#2]{#1}}
\newcommand{\defref}[1]{\hyperref[def:#1]{Definition #1}}


%% Axioms %%
\newtcolorbox[]{axiom}[2][]{%
    rounded corners,
    colframe=xdarkgreen,
    colback=xdarkgreen!15,
    coltitle=white,
    label={ax:#2},
    title={\large \textbf{Axiom of #2}},
    #1
}

%% Exercises %%

\newtcolorbox[]{exercise}[2][]{%
    rounded corners,
    colframe=xdarkgray,
    colback=xgray!15,
    coltitle=white,
    label={ex:#2},
    title={\large \textbf{Exercise: #2}},
    #1,
}

\newcommand{\defex}[1]{\hyperref[ex:#1]{Exercise #1}}

%% Examples %%
\newcounter{ExampleCounter}
\newtcolorbox[use counter=ExampleCounter, number within=chapter]{example}[2][]{%
    rounded corners,
    colframe=teal,
    colback=teal!15,
    coltitle=white,
    label={ex:\theExampleCounter},
    title={\large \textbf{Example \theExampleCounter: #2}},
    #1,
}

%% Corollaries %%
\newtcolorbox[]{corollary}[3][]{%
    rounded corners,
    colframe=xgrayblue,  
    colback=xgrayblue!15,        
    coltitle=white,
    title={\large \textbf{Corollary to #3: #2}},
    #1
}

\newtcolorbox[use counter=BoxCounter, number within=section]{lemma}[1][]{%
    rounded corners,
    colframe=xpurple,
    colback=xpurple!15,
    coltitle=white,
    label={lm:\theBoxCounter},
    title={\large \textbf{Lemma \theBoxCounter}},
    #1,
}

%-%-%-%-%-%-%-%-%-%-%-%-%-%-%-%-%-%-%-%-%-%-%-%-%-%-%-%-%-%-%-%-%-%-%-%-%-%-%

%% MATH PACKAGES, ENVIRONMENTS, COMMANDS %%
%-%-%-%-%-%-%-%-%-%-%-%-%-%-%-%-%-%-%-%-%-%-%-%-%-%-%-%-%-%-%-%-%-%-%-%-%-%-%
%You'll need your own packages, theroem types, and commands.

\usepackage{tikz}
\usetikzlibrary{shapes,backgrounds,calc,patterns}
\usepackage{amsmath,amsthm} 
\usepackage{mathtools}
\usepackage{amssymb}
\usepackage{mdframed}
\usepackage{mathrsfs}
\usepackage{changepage}
\usepackage{multicol}
\usepackage{slashed}
\usepackage{enumerate}
\usepackage{booktabs}
\usepackage{enumitem}
\usepackage{lipsum}  %This package lets us generate random text for example purposes.


%%% Custom Comands %%%
% Natural Numbers 
\newcommand{\N}{\ensuremath{\mathbb{N}}}

% Whole Numbers
\newcommand{\W}{\ensuremath{\mathbb{W}}}

% Integers
\newcommand{\Z}{\ensuremath{\mathbb{Z}}}

% Rational Numbers
\newcommand{\Q}{\ensuremath{\mathbb{Q}}}

% Real Numbers
\newcommand{\R}{\ensuremath{\mathbb{R}}}

% Complex Numbers
\newcommand{\C}{\ensuremath{\mathbb{C}}}

\newcommand{\I}{\ensuremath{\mathbb{I}}}

\newcommand{\pfs}{\noindent\makebox[\linewidth]{\rule{\textwidth}{0.4pt}}\vspace{0.5cm}}

\newcommand{\caution}{\marginnote{\includegraphics[width=2em]{caution.png}}}

\newmdenv[
  topline=false,
  bottomline=false,
  rightline=false,
  leftline=true,
  linewidth=1.5pt,
  linecolor=black, % default color, will be overridden in custom commands
  innertopmargin=0pt,
  innerbottommargin=0pt,
  innerrightmargin=10pt,
  innerleftmargin=10pt,
  leftmargin=0pt,
  rightmargin=0pt,
  skipabove=\topsep,
  skipbelow=\topsep,
]{customframedproof}

\newenvironment{proofpart}[2][black]{
    \begin{mdframed}[
        topline=false,
        bottomline=false,
        rightline=false,
        leftline=true,
        linewidth=1pt,
        linecolor=#1!40, % Custom color
        innertopmargin=10pt,
        innerbottommargin=10pt,
        innerleftmargin=10pt,
        innerrightmargin=10pt,
        leftmargin=0pt,
        rightmargin=0pt,
        skipabove=\topsep,
        skipbelow=\topsep
    ]
    \noindent
    \begin{minipage}[t]{0.08\textwidth} % Adjust width as needed
        \textbf{#2} % The label for the implication (Rightarrow or Leftarrow)
    \end{minipage}%
    \begin{minipage}[t]{0.90\textwidth} % Adjust width as needed
        \begin{adjustwidth}{0pt}{0pt} % No additional indentation
}{
    \end{adjustwidth}
    \end{minipage}
    \end{mdframed}
}


\newcommand{\sbseteqpf}[5][We need to show these expressions are subsets of each other in order to prove they are equivalent.]{%
    \begin{customframedproof}[linecolor=#5]
        \begin{proof}
                #1
            \begin{proofpart}[#5]{(\(\subseteq\))}
                #2
            \end{proofpart}
    \vspace{2mm}
            \begin{proofpart}[#5]{(\(\supseteq\))}
                #3
            \end{proofpart}
    \vspace{2mm}
            #4
        \end{proof}
    \end{customframedproof}
}

% \newcommand{\iffpf}[4]{
%     \begin{customframedproof}[linecolor=#4]
%         \begin{proof}
%             We show this by proving both implications:
%             \begin{proofpart}[#4]{(\(\Rightarrow\))}
%                 #1
%             \end{proofpart}
%     \vspace{2mm}
%             \begin{proofpart}[#4]{(\(\Leftarrow\))}
%                 #2
%             \end{proofpart}
%     \vspace{2mm}
%             #3
%         \end{proof}
%     \end{customframedproof}
% }

\newcommand{\iffpf}[5][We show this by proving both implications:]{
    \begin{customframedproof}[linecolor=#5]
        \begin{proof}
            #1 % This is the optional text, with a default value
            \begin{proofpart}[#5]{(\(\Rightarrow\))}
                #2
            \end{proofpart}
    \vspace{2mm}
            \begin{proofpart}[#5]{(\(\Leftarrow\))}
                #3
            \end{proofpart}
    \vspace{2mm}
            #4
        \end{proof}
    \end{customframedproof}
}



% Command for lemma proofs with a predefined color
\newcommand{\expf}[1]{
    \begin{customframedproof}[linecolor=xgray]
        \begin{proof}
        #1
        \end{proof}
    \end{customframedproof}
}


% Command for theorems proofs with a predefined color
\newcommand{\tpf}[1]{
    \begin{customframedproof}[linecolor=xblue]
        \begin{proof}
        #1
        \end{proof}
    \end{customframedproof}
}
% Command for corollaries proofs with a predefined color
\newcommand{\cpf}[1]{
    \begin{customframedproof}[linecolor=purplehdx]
        \begin{proof}
        #1
        \end{proof}
    \end{customframedproof}
}

\newcommand{\sol}[1]{
    \begin{customframedproof}[linecolor=xgray]
        \begin{solution}
        #1
        \end{solution}
    \end{customframedproof}
}

\newcommand{\lesol}[1]{
    \begin{customframedproof}[linecolor=teal]
        \begin{solution}
        #1
        \end{solution}
    \end{customframedproof}
}

\newcommand{\pb}[2]{
    \begin{exercise}{#1}[label=ex:#1]
    #2
    \end{exercise}
}

\def \proofDistance {10pt}
\newenvironment{solution}
  {\textit{Solution.}}

% Command for theorems proofs with a predefined color
\newcommand{\epf}[1]{
    \begin{customframedproof}[linecolor=teal]
        \begin{proof}
        #1
        \end{proof}
    \end{customframedproof}
}

\newcommand{\disabs}[1]{\ensuremath{\displaystyle\left|#1\right|}}
\newcommand{\limn}[1]{\ensuremath{\displaystyle \lim_{n \rightarrow \infty}}#1}

\renewcommand{\theenumi}{\alph{enumi}} 
\renewcommand{\labelenumi}{(\theenumi)}


\usepackage[toc, hyperfirst]{glossaries}

\makeglossaries

\newglossaryentry{The Dirichlet Function}
{
name=The Dirichlet Function,
description={\(f(x)=
\begin{cases}
    x = 1 &\text{ if } x \in \Q \\
    x = 0 &\text{ if } x \notin \Q
\end{cases}\)}
}

\newglossaryentry{least upper bound}
{   
    name=least upper bound,
    description={A number \(r\) is a least upper bound of \(A\) if \(r\) is an upper bound, for all upper bounds, \(s\), \(r \leq S\)}
}

\newglossaryentry{upper bound}
{
    name=upper bound,
    description={For a number \(s \in \R\) such that for every \(x \in A\), \(x \leq s\)}
}

\newglossaryentry{bounded}
{
    name=bounded,
    description={For a set \(A \subseteq \R\) there exists \(M > 0\) such that for all \(x \in A\), \(|x| \leq M\)}
}

\newglossaryentry{natural numbers}
{
    name=natural numbers,
    description={Contains all positive, non-zero, and non-fractional numbers. Expressed as \(\N = \{1,2,3,4\dots\}\)}
}

\newglossaryentry{integers}
{
    name=integers,
    description={Contains all non-fractional numbers. Expressed as: \(\Z = \{\dots-2,-1,0,1,2\dots\}\)}
}

\newglossaryentry{rational numbers}
{
    name=rational numbers,
    description={Contains all numbers, except irrational numbers. Expressed as: \(\Q = \{\frac{p}{q} \ | \ p,q \in \Z, q \ne 0\}\)}
}

\newglossaryentry{triangle inequality}
{
    name=triangle inequality,
    description={For any \(a,b \in \R\), \(|a + b| \leq  |a| + |b|\)}
}

\newglossaryentry{pre-image}
{
    name=pre-image,
    description={\(f \colon A \rightarrow B\). If \(E \subseteq B\), \(f^{-1}(E) = \{x\in A \ | \ f(x) \in E\}\). With \(g(x) = 2x + 1\) from image}
}

\newglossaryentry{image}
{
    name=image,
    description={\(f\colon A \rightarrow B\). Let \(E \subseteq A\) such that \(f(E) = \{f(x) \ | \ x \in E\}\) \(g(x) = 2x + 1\) where \(g\) is defined as \(g \colon \R \rightarrow \R\).}
}

\newglossaryentry{De Morgan's Laws}
{
    name=De Morgan's Laws,
    description={\((A \cup B)^c = A^c \cap B^c\) and \((A \cap B)^c = A^c \cup B^c\)}
}

\newglossaryentry{one-to-one}
{
    name=one-to-one,
    description={If \(a_1 \ne a_2\) in \(A\) implies that \(f(a_1) \ne f(a_2)\) in \(B\)}
}

\newglossaryentry{onto}
{
    name=onto,
    description={Given any \(b\in B\), it is possible to find an element \(a \in A\) for which \(f(a) = b\).}
}

\newglossaryentry{maximum value}
{
    name=maximum value,
    description={If \(m \in A\) and for all \(x\in A\), \(x \leq m\)}
}



% end of preamble
%-%-%-%-%-%-%-%-%-%-%-%-%-%-%-%-%-%-%-%-%-%-%-%-%-%-%-%-%-%-%-%-%-%-%-%-%-%-%

\begin{document}

\numberwithin{BoxCounter}{section}

	
%-%-%-%-%-%-%-%-%-%-%-%-%-%-%-%-%-%-%-%-%-%-%-%-%-%-%-%-%-%-%-%-%-%-%-%-%-%-%
 %%% COVER PAGE %%%
 %-%-%-%-%-%-%-%-%-%-%-%-%-%-%-%-%-%-%-%-%-%-%-%-%-%-%-%-%-%-%-%-%-%-%-%-%-%-%

% Do not use all caps.
\newcommand{\titlestandin}[0]{Real Analysis}
\newcommand{\cussubtitle}[0]{MATH 350}
% Date format should be like: "January 1, 2001"
\newcommand{\startdate}[0]{August 26, 2024}
\newcommand{\customenddate}[0]{December 2, 2024}
% Be sure to include degree recognition (e.g., B.S., M.S., Ph.D)
\newcommand{\professor}[0]{Prof. Christopher Camfield, Ph.D.}

%-%-%-%-%-%-%-%-%-%-%-%-%-%-%-%-%-%-%-%-%-%-%-%-%-%-%-%-%



\begin{titlepage}
\begin{center}

\vspace*{-2cm}
\includegraphics[width=0.8\textwidth]{images/Hendrix Logo.png}\\
\vfill

% Horizontal line above the title in 'horange' color
\textcolor{horange}{\rule{\textwidth}{1.0pt}}

\vspace{2em}

{\huge \textbf{\titlestandin}}

\vspace{1em} % Space between the title and the bottom line

\textcolor{horange}{\rule{\textwidth}{1.0pt}}

\vspace*{1\baselineskip}

{\LARGE \textbf{\cussubtitle}}

\begin{large}
\vspace*{2\baselineskip}

\textit{Start}  \\[1ex]
{\scshape \startdate} \\[0.3\baselineskip] % Year published

\vspace*{1\baselineskip}

\emph{Author} \\[1ex]
%Submitted by \\[\baselineskip]
{\Large Paul Beggs \\ \par} % Editor list
{\href{mailto:BeggsPA@Hendrix.edu}{{BeggsPA@Hendrix.edu}}}\\ % Editor affiliation

\vspace*{1\baselineskip}

\textit{Instructor} \\[1ex] % Tagline(s) or further description
\professor

\vspace*{1\baselineskip}

\textit{End}\\[1ex]
{\scshape  \customenddate} \\[0.3\baselineskip] % Year published

\thispagestyle{empty} 

\end{large}
\end{center}
\end{titlepage}
\pagebreak


%-%-%-%-%-%-%-%-%-%-%-%-%-%-%-%-%-%-%-%-%-%-%-%-%-%-%-%-%-%-%-%-%-%-%-%-%-%-%

%%% Table of Contents %%%

% No need to edit.

%-%-%-%-%-%-%-%-%-%-%-%-%-%-%-%-%-%-%-%-%-%-%-%-%-%-%-%-%-%-%-%-%-%-%-%-%-%-%
	
	\thispagestyle{plain}
	% Table of Contents will automatically update based on included chapters
	% You may need to compile twice for the toc to update
	\begin{spacing}{1} %Can change spacing between entries in toc
		\renewcommand{\contentsname}{\Large\textbf{Table of Contents}} %Can rename here
		\tableofcontents
		\addtocontents{toc}{\protect\enlargethispage{\baselineskip}} 
	\end{spacing}


%-%-%-%-%-%-%-%-%-%-%-%-%-%-%-%-%-%-%-%-%-%-%-%-%-%-%-%-%-%-%-%-%-%-%-%-%-%-%



%-%-%-%-%-%-%-%-%-%-%-%-%-%-%-%-%-%-%-%-%-%-%-%-%-%-%-%-%-%-%-%-%-%-%-%-%-%-%



%%% CHAPTERS %%%

% The document will automatically update (including numbering and toc)
% You may need to compile twice to update the toc
% Witness magic!
	
% Chapter files -- containing the actual text of the document
 % These are in the chapters folder
% All images (for any chapter) must be placed in the images folder

% Example chapters are included with examples on formatting boxes and figures


%-%-%-%-%-%-%-%-%-%-%-%-%-%-%-%-%-%-%-%-%-%-%-%-%-%-%-%-%-%-%-%-%-%-%-%-%-%-%

% \chapter{The Real Numbers}\label{Real nums}
% \vspace*{-0.25in}
% \section{Types of Numbers}

The \gls{natural numbers}, $\N$:
\begin{itemize}
    \item No additive inverse.
    \item You can:
    \begin{itemize}
        \item Add,
        \item Multiply
    \end{itemize}
\end{itemize}

The \gls{integers}, $\Z$ are known as a Group (more specifically, a ``ring'').

\begin{itemize}
    \item You can:
    \begin{itemize}
        \item Add,
        \item Multiply,
        \item Subtract
    \end{itemize}
\end{itemize}

The \gls{rational numbers}, $\Q$ are known as a ``Field.''

\begin{itemize}
    \item You can:
    \begin{itemize}
        \item Add,
        \item Subtract,
        \item Multiply,
        \item Divide
    \end{itemize}
\end{itemize}

A problem that rational numbers could not explain: The $45, 45, 90$ triangle had a hypotenuse of $\sqrt{2}$. This did not exist at the time, so it was simply $c^2 = 2$. Therefore, new numbers needed to be invented.

\begin{theorem}
    There does not exist a rational number $r$ such that $r^2 = 2$.
\end{theorem}

\tpf{
    Suppose there exists a rational number $r$ such that $r^2 = 2$. Since $r$ is rational, there exists $p,q \in \Z$ such that $r = \frac{p}{q}$. We can assume the $p$ and $q$ have no common factors. (If not, we can factor out the common factor.) By our assumption, 
    \begin{align*}
        r^2 &= 2 \\
        \frac{p^2}{q^2} &= 2 
    \end{align*}
    It follows that, \begin{align*}
        p^2 = 2q^2
    \end{align*}
    Such that $p^2$ is an even number because if $p$ were odd, then $p^2$ would be odd. There exists $x \in \Z$ such that $p = 2x$. Recall that $p^2 = 2q^2$. Thus\begin{align*}
        (2x)^2 &= 2q^2 \\
        4x^2 &= 2q^2 \\
        2x^2 &= q
    \end{align*} Thus, $q^2$ is even. Hence $q$ is also even. So $p$ and $q$ are both divisible by 2. This contradicts that $p$ and $q$ have no common factors. Thus, our supposition is false. Therefore, there does not exist a rational number $r$ such that $r^2 = 2$ 
}

So we are going to work with a larger set called the real numbers, $\R$. 
\begin{itemize}
    \item $\N \subset \Z \subset \Q \subset \R$
    \item You can:
    \begin{itemize}
        \item Add,
        \item Subtract,
        \item Multiply,
        \item Divide
    \end{itemize}
    \item In other words, all field axioms apply.
    \item Totally ordered set for any $x,y \in \R$. Thus, one of these are true:
    \begin{enumerate}
        \item $x < y$,
        \item $x > y$,
        \item $x = y$
    \end{enumerate}
    \item  Think of it as a number line.
    \item $\Q$ is dense: \\
    If $a,b \in \Q$ with $a \ne b$, there exists $c \in \Q$ which is between $a$ and $b$ such that $a < c < b$. One example is $\frac{a + b}{2}$.
    \item $\Q$ is not \textit{complete}, but $\R$ is.
    \begin{itemize}
        \item \textit{Complete:} Think, ``no gaps.''
    \end{itemize}
\end{itemize}

\section{Preliminaries}

Things to remember from Intro and Discrete.

\begin{table}[htbp]
\begin{tabular}{c||c}
Set Notation & Complement \\
\hline 
$x \in A$ & $A^c$ (not $\overline{A}$) \\
$A \cup B$ & $\R \setminus A$ \\
$A \cap B$
\end{tabular}
\end{table}

\begin{itemize}
\item $\displaystyle \bigcup^\infty_{n=1} A_n = A_1 \cup A_2 \cup A_3 \cup \dots$. \\

\item $\displaystyle \bigcap^\infty_{n=1} A_n = A_1 \cap A_2 \cap \dots$. \\

\item \textbf{\Gls{De Morgan's Laws}} \\
\end{itemize}


\pfs

\subsection{Infinite Unions and Intersections} \hfill

For each $n \in \N$, define $A_n = \{n, n+1, n+2,\dots\} = \{k \in \N \ | \ k \geq n\}$. In other words, each subsequent element in the subset will start at $n$. For example, $A_1 = \{1,2,\dots\}$, whereas $A_5 = \{5,6,\dots\}.$   \\

$\bigcup^\infty_{n=1} A_n = \N$. To show a number $\in \N$ belongs in the set $A_n$, we can start with that, $k \in \N$. Then $k \in A_k$. Thus, $k \in A_k \subseteq \bigcup^\infty_{n=1}A_n$. therefore, $\N \subseteq \bigcup^\infty_{n=1}A_n$. \\

$\bigcap^\infty_{n=1}A_n = \emptyset$. Obviously, we know that the empty set is a subset of $A_n$, but to prove that $\bigcap^\infty_{n=1}A_n$ is a subset of the empty set, we should suppose a $k \in \N$ such that $k \in \bigcap^\infty_{n=1}A_n$. Notice that $k \notin \bigcap^\infty_{n=1}A_n$. So, $\bigcap^\infty_{n=1}A_n = \emptyset$.

\pfs

\subsection{Functions and Notation} \hfill

$f \colon A \rightarrow B$ where $f$ is a function, $A$ is a domain, and $B$ is the co-domain. Thus, $f(x) = y$ such that $x \in A$ and $y \in B$. \\

Some definitions to keep in mind: \\

\textbf{\Gls{The Dirichlet Function}} \\

\textit{[Refer to notepaper for these following definitions]} \\

\textbf{\Gls{image}} \\

Example: $g \ \colon \ \R \rightarrow \R$, when we say $y \in g(A)$ implies $\exists x$ such that $g(x) = y$

\textbf{\Gls{pre-image}} \\

\textbf{\Gls{triangle inequality}}: \\
The most common application: For any $a,b,c \in \R$, $|a - b| \leq |a-c| + |c-b|$, with the intermediate step of $a - b = (a-c) + (c-b)$.

\subsection{Common Strategies for Analysis Proofs} \hfill

\setcounter{BoxCounter}{5}
\begin{theorem}
    Let $a,b \in \R$. Then, $$a = b \text{ if and only if for all } \epsilon > 0, |a-b| < \epsilon.$$
\end{theorem}

\iffpf
    {Assume $a = b$. Let $\epsilon > 0$. Then $|a-b| = 0 < \epsilon$}
    {Assume for all $\epsilon > 0$, $|a-b| < \epsilon$. Suppose $a \ne b$. Then $a - b \ne 0$. So, $|a-b| \ne 0$. Now, Consider $\epsilon_0 = |a-b|$. By our assumption we know that $|a-b| < \epsilon_0$. It is not true that $|a-b| < |a - b|$. Therefore, it must be the case that $a = b$.}
    {Therefore, by showing both sides of the implication accomplish the same thing as the other side, we know that $a = b \text{ if and only if for all } \epsilon > 0, |a-b| < \epsilon.$}
    {xblue}

\subsection{Mathematical Induction} \hfill

\textit{Inductive Hypothesis:} Let $x_1 = 1$. For all $n \in \N$, let $x_{n+1} = \frac{1}{2}x_n + 1$. \\

\textit{Inductive Step:} $x_1 = 1, x_2 = 1.5, x_3 = 1.75, x_4 = 1.875$. \\

\begin{example}
    {Induction} The sequence $(x_n)$ is increasing. In other words, for all $n\in \N$, $x_n \leq x_{n+1}$. 
\end{example}

\epf{
    Suppose the sequence $(x_n)$ is increasing. We will prove this point by using induction. \\
    
    \textbf{Base Case:} We see that $x_1 = 1$ and $x_2 = 1.5$. Thus, $x_1 \leq x_2$. \\
    
    \textbf{Inductive Hypothesis:} For $n \in \N$, assume $x_n \leq x_{n+1}$. \\
    
    \indent \textit{Scratch work:} We want: $x_{n+1} \leq x_{n+2}$. We know: $x_{n+1} = \frac{1}{2}x_{n+1}+1$. \\
    
    \textbf{Inductive Step:} Then $\frac{1}{2}x_n \leq \frac{1}{2}x_{n+1}$. Hence, $\frac{1}{2}x_{n} + 1 \leq \frac{1}{2}x_{n+1} + 1$. Therefore we have proven through induction that, $x_{n+1} \leq x_{n+2}$.
}

\subsection*{Exercises}

\begin{exercise}
    {1.2.3}Decide which of the following represent true statements about
the nature of sets. For any that are false, provide a specific example where the
statement in question does not hold.
\begin{enumerate}
    \item If $A_1 \supseteq A_2 \supseteq A_3 \supseteq A_4\dots$ are all sets containing an infinite number of
elements, then the intersection $\bigcap^\infty_{n=1} A_n$ is infinite as well.
    \item If $A_1 \supseteq A_2 \supseteq A_3 \supseteq A_4\dots$ are all finite, nonempty sets of real numbers,
then the intersection $\bigcap^\infty_{n=1} A_n$ is finite and nonempty.
    \item $A \cap (B \cup C) = (A \cap B) \cup C$
    \item $A \cap (B \cap C) = (A \cap B) \cap C$
    \item $A \cap (B \cup C) = (A \cap B) \cup (A \cap C)$
\end{enumerate}
\end{exercise}

\sol{\hfill
    \begin{enumerate}
        \item This is false. Consider the following as a counterexample: If we define $A_1$ as $A_n = \{n, n+1, n+2,\dots\} = \{k \in \N \ | \ k \geq n\}$, we can see why the intersection of these sets of infinite numbers are actually empty. Consider a number $m$ that actually satisfies $m \in A_n$ for every $A_n$ in our collection of sets. Because $m$ is not an element of $A_{m + 1}$, no such $m$ exists and the intersection is empty.
        \item This is false. Consider the following as a counter example: Since $A_n$ is finite, consider the following sets: $A_1 = \{1\}$, $A_2 = \{2\}$, and so on. The intersection of these sets is empty because none of the sets contain numbers that are present in the other sets.
        \item This is false. Let the following sets be defined as $A = \{1\}$, $B = \{2\}$, $C = \{3\}$. If we start on the left side of the equation: $A \cap (B \cup C)$ implies $\{1\} \cap (\{2\} \cup \{3\})$ implies $\{1\} \cap \{2, 3\}$ implies $\emptyset$. From the right: $(A \cap B) \cup C$ implies $ (\{1\} \cap \{2\}) \cup \{3\}$ implies $ \emptyset \cup \{3\}$ implies $ \{\emptyset, 3\}$.
        \item  This is true. Let $x \in A$, $y \in B$, and $z \in C$. Consider the following:
        \begin{enumerate}[label=(\arabic*)]
            \item \textbf{From the left of the equation:} Inside the parenthesis, either the sets share an element and combine, or they do not share an element and are empty. This means we have two possibilities:
            \begin{enumerate}[label=\roman*.]
                \item \textbf{If the set inside the parenthesis is not empty}, then we have $x$ and $y \in A \cap B$. Then it could be the case that $x$ and $y \in C$. From this, we know 3 things can occur. Either every set contains only $y$, or $x$ and $y$.
                \item \textbf{If the set inside the parenthesis is empty}, then the resulting intersection with $C$ would be the empty set.
            \end{enumerate}
            \item \textbf{From the right of the equation:} Similar to (1), either the sets share an element and combine, and they do not share an element and are empty. This means we have two possibilities:
            \begin{enumerate}[label=\roman*.]
                \item \textbf{If the set inside the parenthesis is not empty}, then we have $y$ and $z \in B \cap C$. Then it could be the case that $y$ and $z \in A$. From this, we know 2 things can occur. Either every set contains only $y$, or $y$ and $z$.
                \item \textbf{If the set inside the parenthesis is empty}, then the resulting intersection with $A$ would be the empty set.
            \end{enumerate}
            \item In every case, either we get the empty set, or end up in a situation where every set contains just $y$ or all elements. Hence, the two sides of the equation will always equal each other. 
        \end{enumerate}
        \item This is true. Because $B \cap C$ contains all the elements in $B$ or $C$, so when the intersection is taken on either side of the equation, you will have the same elements regardless. (This statement is much shorter than the explanation given in the previous problem because I am unsure if I even a proof at all. I'm testing both waters here.) 
    \end{enumerate}
}

\newpage
\begin{exercise}
    {1.2.5}\Gls{De Morgan's Laws}. Let $A$ and $B$ be subsets of $\R$.
    \begin{enumerate}
        \item  If $x \in (A \cap B)^c$, explain why $x \in A^c \cup B^c$. This shows that $(A \cap B)^c \subseteq A^c \cup B^c$.
        \item Prove the reverse inclusion $(A \cap B)^c \supseteq A^c \cup B^c$, and conclude that $(A \cap B)^c = A^c \cup B^c$.
        \item Show $(A \cup B)^c = A^c \cap B^c$ by demonstrating inclusion both ways.
    \end{enumerate}
\end{exercise}

\sol{\hfill
    \begin{enumerate}
        \item If $x \in (A \cap B)^c$, and we know that $A^c = \{x \in \R \ \colon x \notin A\}$, then we know $x$ must cannot exist in $A^c$ and $B^c$ because $(A \cap B)^c = \{x \in \R \ \colon x \notin (A \cap B)\}$. Thus, $x$ is in either $A^c$ or $B^c$. Put another way $x \in A^c \cup B^c$. Since we have shown that an element that started in $(A \cap B)^c$ ended up in $A^c \cup B^c$, then we know $(A \cap B)^c \subseteq A^c \cup B^c$.
        \item Assume there exists a $y \in A^c \cup B^c$. Thus, it must be the case that $y \notin A$ or $y \notin B$. Hence, $y$ cannot be exist in both sets at the same time, so $y \in (A \cap B)^c$. Because we have taken an element that started in $A^c \cup B^c$ and have shown that it exists in $(A \cap B)^c$, we have proven $A^c \cup B^c \subseteq (A \cap B)^c$. 
        \item\hfill \sbseteqpf
    {Assume there exists $x \in (A \cap B)^c$, and we know that $A^c = \{x \in \R \ \colon x \notin A\}$, then we know $x$ must cannot exist in $A^c$ and $B^c$ because $(A \cap B)^c = \{x \in \R \ \colon x \notin (A \cap B)\}$. Thus, $x$ is in either $A^c$ or $B^c$. Put another way $x \in A^c \cup B^c$. Since we have shown that an element that started in $(A \cap B)^c$ ended up in $A^c \cup B^c$, then we know $(A \cap B)^c \subseteq A^c \cup B^c$.}
    {Now assume there exists a $y \in A^c \cup B^c$. Thus, it must be the case that $y \notin A$ or $y \notin B$. Hence, $y$ cannot be exist in both sets at the same time, so $y \in (A \cap B)^c$. Because we have taken an element that started in $A^c \cup B^c$ and have shown that it exists in $(A \cap B)^c$, we have proven $A^c \cup B^c \subseteq (A \cap B)^c$.}
    {Therefore, we have shown through proving both sides of the implication, that these two statements are logically equivalent. In that, all elements of $A^c \cup B^c$ are the same elements that are in $(A \cap B)^c$}
    {xgray}
    \end{enumerate}
}



\begin{exercise}
    {1.2.7}Given a function $f$ and a subset $A$ of its domain, let $f(A)$
represent the range of $f$ over the set $A$; that is, $f(a) = \{f(x) \colon x \in A\}$.
\begin{enumerate}
    \item Let $f(x) = x^2$. if $A = [0,2]$ (the closed interval $\{x \in \R \colon 0 \leq x \leq 2\}$) and $B = [1,4]$, find $f(A)$ and $f(B)$. Does $f(A\cap B) = f(A) \cap f(B)$ in this case? Does $f(A \cup B) = f(A) \cup f(B)$?
    \item Find two sets $A$ and $B$ for which $f(A\cap B) \ne f(A) \cap f(B)$.
    \item Show that, for an arbitrary function $g \ \colon \ \R \rightarrow \R$, it is always true that $g(A \cap B) \subseteq g(A) \cap g(B)$ for all sets $A,B\subseteq \R$.
    \item Form and prove a conjecture about the relationship between $g(A \cup B)$ and $g(A) \cup g(B)$ for an arbitrary function $g$.
\end{enumerate}
\end{exercise}

\sol{\hfill
    \begin{enumerate}
        \item Since $f(x) = x^2$, the intervals of $f(A)$ would be $[0,4]$ and $f(B)$ would be $[1,16]$. The interval of the intersection of $A \cap B$ is $[1,2]$. Take this through our function, we get $f(A \cap B) = [1,4]$. On the other side of the equation, we already know the intervals of $f(A)$ and $f(B)$, and the intersection of theirs would be $[1,4]$. So they do equal each other. We know $f(A \cup B)$ and $f(A) \cup f(B)$ will be equivalent because $f(A \cup B)$ has an interval of $[0,16]$, and $f(A) \cup f(B)$ also has an interval of $[0,16]$ because taking the union of $[0,4] \cup [1,16]$ is $[0,16]$.
        \item Two sets could be $A = [5,6]$ and $B = [0,0]$. Because the sets have nothing in common even after taking their function, they do not equal each other. 
        \item \hfill \expf{Let $x \in g(A \cap B)$. Using the definition of function, we know there exists a $y \in A \cap B$ to which that $y$ is mapped to as $g(y) = x$. From the definition of intersection, we know $y \in A$ and $y \in B$ such that $x = g(y) \in g(A)$ and $x = g(y) \in g(B)$ because $y \in A \cap B$. Putting it together, we have $x \in g(A) \cap g(B)$ thus proving $g(A \cap B) \subseteq g(A) \cap g(B)$}
        \item Conjecture: For any function $g$ defined as $g \ \colon \R \rightarrow \R$ and for any subsets $A,B \subseteq \R$, the following holds: $$g(A \cup B) = g(A) \cup g(B)$$ 
        \sbseteqpf
        {Take any element $x \in g(A \cap B)$. By definition of function, we know there exists some $y \in A \cup B$ such that $g(x) = y$. From the definition of union, we know $y \in A$ or $y \in B$ such that $x = g(y) \in g(A)$ or $x = g(y) \in g(B)$ or both. Putting it together, we have $x \in g(A) \cup g(B)$ thus proving $g(A \cup B) \subseteq g(A) \cup g(B)$.}
        {Take any element $p \in g(A) \cap g(B)$. By definition of union, we know $p$ is either in $g(A)$ or $g(B)$ or both. From the definition of function, we know that if $p \in g(A)$ or $p \in g(B)$ then there exists some $q \in A$ or $q \in B$ such that $g(q) = p$. Putting it together, we have $q \in A \cup B$. Moreover, this means $p = g(x) \in g(A \cup B)$. And since $p \in g(A) \cup g(B)$ implies $p \in g(A \cup B)$, we know $g(A) \cup g(B) \subseteq g(A \cup B)$.}
        {Therefore, since we have proven that both expressions are functions of each other, we have proved that they are equal.}
        {xgray}
    \end{enumerate}
}

\begin{exercise}
    {1.2.8}Given a function $f \ \colon A \rightarrow B$ can be defined as either \gls{one-to-one} or \gls{onto}, give an example of each or state that the request is impossible:
    \begin{enumerate}
        \item $f \ \colon \N \rightarrow \N$ that is 1-1 but not onto.
        \item $f \ \colon \N \rightarrow \N$ that is onto but not 1-1.
        \item $f \ \colon \N \rightarrow \Z$ that is 1-1 and onto.
    \end{enumerate}
\end{exercise}

\sol{\hfill
    \begin{enumerate}
        \item The function $f(a) + 1$ is 1-1 because when \begin{align*}
            f(a_1) &= f(a_2)\\
            a_1 + 1 &= a_2 + 1 \\
            a_1 &= a_2
        \end{align*}
        However, the function is not onto because the entire co-domain is not covered. That being 1.
        \item \textit{We need to find a function that will cover every entry in the co-domain, while also avoiding a scenario where $a_1 = a_2$...} Consider the function, 
        $$f(a) =
        \begin{cases}
            a &\text{ if \textit{a} is odd,}\\
            a - 1 &\text{ if \textit{a} is even}
        \end{cases}$$
        This function is onto because every natural number is covered, but it is not 1-1 because $a_1 \ne a_2 - 1$.
        \item This request is not possible. There is no way to map every natural number to every integer because we are simply missing 0! (Not 0 factorial, we do have the number 1, I just mean the number 0 in a exclamatory sense.)
    \end{enumerate}
}
\setcounter{section}{2}
\section{Axiom of Completeness}

Think about $\Q$ and $\R$.
    \begin{itemize}
        \item Both are fields.
        \begin{itemize}
            \item Both have $+,-,\times,\div$ operations.
        \end{itemize}
        \item Both are totally ordered
        \begin{itemize}
            \item $a < b$,
            \item $a > b$, 
            \item or $a = b$
        \end{itemize}
        \item $\R$ is complete. $\Q$ is not.
    \end{itemize}
    
\begin{axiom}
    {Completeness}Every nonempty set of real numbers that is \textbf{\gls{bounded}} has a \textbf{\gls{least upper bound}}.
\end{axiom}

\textbf{Note that \glspl{upper bound} are not unique!} For example, consider the line, $A$, from $0$ to $1$. There are infinitely many upper bounds past $1$ because $A$ is bounded. \\

We often call the least upper bound the \textit{supremum} of a set. Example: \\

Imagine a number line from $(1,8)$. Note that parenthesis mean $<$ and not $\leq$. Hence, the supremum is $8$. Wrote simply as $\sup A$. \\

\begin{example}
    {Supremum}Consider a set, $B = [-5,-2] \cup (3,6) \cup \{13\}$. What is the supremum?
\end{example}

\sol{
    $\sup B = 13$
}

At the other end of the set, we have the following:
\begin{itemize}
    \item lower bounds,
    \item greatest lower bound
    \item often called infimum.
\end{itemize}

The infimum of the previous example would be $\inf B = -5$.

\begin{example}
    {}Consider the set, $\C = \{\frac{1}{n} \ \colon \ n \in \N\}$. What is the supremum and the infimum?
\end{example}

\lesol{
    $\sup \C = 1$, $\inf \C = 0$.
}

\begin{example}
    Let $A \subseteq \R$ be nonempty and bounded above, and let $c \in \R$. Define the set $c + A$ by $$c + A = \{c+a \ \colon \ a \in A\}$$ Then $\sup(c + A) = c + \sup A$.
\end{example}

\lesol{
To properly verify this we focus separately on each part of Definition 1.3.2. Setting $s = \sup A$, we see that $a \leq s$ for all $a \in A$, which implies $c + a \leq c + s$ for all $a \in A$. Thus, $c + s$ is an upper bound for $c + A$ and condition (i) is verified. For (ii), let $b$ be an arbitrary upper bound for $c + A$; i.e., $c + a \leq b$ for all $a \in A$. This is equivalent to $a \leq b - c$ for all $a \in A$, from which we conclude that $b - c$ is an upper bound for $A$. Because $s$ is the least upper bound of $A$, $s \leq b - c$, which can be rewritten as $c + s \leq b$. This verifies part (ii) of Definition 1.3.2, and we conclude $\sup (c + A) = c + \sup A$.
}

Why do we need to include infimum and supremum? Don't we have the max and min of a set already? Well, what exactly do we mean by the \gls{maximum value} of a set? \\

We say $m \in \R$ is the \textit{maximum} of $A$ if $m \in A$ and for all $x\in A$, $x \leq m$. Note that some sets have a maximum and some sets do not. You cannot refer to a maximum without first knowing it exists. This is the same with minimums. Suprema do not exist if the set is bounded.

\begin{lemma}
    {1.3.8}Assume $s$ is an \gls{upper bound} for a set $A \subseteq \R$. Then, $s$ is the supremum of $A$ if and only if for every $\epsilon > 0$, there exists $x\in A$ such that $s - \epsilon < x$.
\end{lemma}

This lemma allows us to take any positive number and take a ``step back.'' In essence, you can verify something as an upper bound if you continuously back up over and over until you cannot back up any longer.

\iffpf
{
Assume $s = \sup A$. Let $\epsilon > 0$. Suppose there are no elements $x$ of $A$ such that $s - \epsilon < x$. Then $s - \epsilon$ would be an upper bound. This contradicts that $s$ is the least upper bound. Therefore, there must exist an element $x \in A$ such that $s - \epsilon < x$. 
}{
Assume for every $\epsilon > 0$, there exists $x \in A$ such that $s - \epsilon < x$. Let $t$ be an upper bound of $A$. Suppose $t < s$. Consider $\epsilon_0 = s - t > 0$. By our assumption, there exists $x \in A$ such that $s - \epsilon_0 < x$. So, $t < x$. This contradicts that $t$ is an upper bound of $A$. So, $t \geq s$. Thus, $s$ is the least upper bound
}{
Therefore, by proving both the right and left implication, we have shown the statement to be true.
}{xpurple}
Analogous statement about infimums: Assume $z$ is a lower bound of a set $A \subseteq \R$. Then $z = \inf A \iff $ for all $ \epsilon >0$, there exists $y \in A$ such that $y < z + \epsilon$.

\subsection*{Exercises}

\begin{exercise}
    {1.3.4}Let $A_1,A_2,A_3\dots$ be a collection of nonempty sets each of which is bounded above. 
    \begin{enumerate}
        \item Find a formula for $\sup(A_1 \cup A_2)$. Extend this to $\sup (\bigcup^n_{k=1}A_k)$.
        \item Consider $\sup(\bigcup^\infty_{k=1} A_k)$. Does the formula in (a) extend to the infinite case?
    \end{enumerate}
\end{exercise}

\sol{\hfill
\begin{enumerate}
    \item Let $A_1$ and $A_2$ be nonempty sets, each bounded above. To find the largest of the two suprema, we can use the following: $\sup(A_1 \cap A_2) = \max\{\sup A_1, \sup A_2\}$. If we extend this notion to $\sup (\bigcup^n_{k=1}A_k)$, we can use the same idea from before and write it as $\sup (\bigcup^n_{k=1}A_k) = \max\{\sup A_1, \sup A_2, \dots, \sup A_n\}$.
    \item The formula does not extend to the infinite case. Consider the counterexample $\bigcup^\infty_{k=1} A_k$ where $A_k := [k, k+1]$. Even though these sets are bounded above, when we take the union of them, we approach infinity, which is not bounded: $\bigcup^\infty_{k=1} A_k = [1,2] \cup [2,3] \cup \dots = [1, \infty)$.
\end{enumerate}
}

\begin{exercise}
    {1.3.5}As in Example 1.3.7, let $A \subseteq \R$ be nonempty and bounded above, and let $c \in \R$. This time define the set $cA = \{ca \ \colon a \in A\}$.
    \begin{enumerate}
        \item If $c \geq 0$, show that $\sup(cA) = c\sup A$.
        \item Postulate a similar type of statement for $\sup(cA)$ for the case $c < 0$.
    \end{enumerate}
\end{exercise}

\sol{
\begin{enumerate}
    \item Let $A \subseteq \R$ be nonempty and bounded above. Define the set $cA := \{ca \ \colon a \in A\}$. From the axiom of completeness, because $A$ is bounded above, we know there is a least upper bound, $s = \sup A$. Following from Example 1.3.7, we see that $a \leq s$ for all $a \in A$ which implies $ca \leq cs$ for all $a \in A$. Thus, $cs$ is an upper bound for $cA$, and the first condition of Definition 1.3.2 is satisfied. For the second condition, we need to look at both $c = 0$ and $c > 0$ to avoid dividing by zero. So, we have two cases:
    \begin{itemize}
        \item \textbf{$c = 0$:} If $c = 0$, then $cA = \{0 \colon \ a \in A\} = \{0\}$. Since the only element in $cA$ is $0$, $\sup(cA) = 0$. Similarly, because $c = 0$, $c \sup A = 0 \cdot \sup A = 0$. Therefore, $\sup(cA) = c \sup(A)$.
        \item \textbf{$c > 0$:} Let $b$ be an arbitrary upper bound for $cA$ and $c > 0$. In other words, $ca \leq b$ for all $a \in A$. This is equivalent to $a \leq b / c$ where $c \ne 0$, from which we can see that $b / c$ is an upper bound for $A$. Because $s$ is the least upper bound of $A$, $s \leq b / c$, which can be rewritten as $cs \leq b$. This verifies the second part of Definition 1.3.2, and we conclude $\sup(cA) = c \sup A$.
    \end{itemize}
    \item Postulate: If $c < 0$, then $\sup(cA) = c \inf (A)$. \\
    % \expf{Assume $c < 0$, Multiplying $a \leq s$ by $c$ reverses the inequality, so $ca \geq cs$ for all $a \in A$. Therefore, $cs$ is a lower bound of for $cA$. To verify this as the lowest upper bound, let $m  = \inf A$. Using the same logic as before, for all $a \in A$, $m \leq a$, so $cm \geq ca$. This shows that $cm$ is the least upper bound of $cA$, and thus $\sup(cA) = cm = c \inf A$.}
\end{enumerate}
}

\begin{exercise}
    {1.3.8}Compute, without proofs, the suprema and infima (if they exist) of the following sets:
\begin{enumerate}
    \item $\left\{\frac{m}{n} : m, n \in \mathbb{N} \text{ with } m < n\right\}$.
    \item $\left\{\frac{(-1)^m}{n} : m, n \in \mathbb{N}\right\}$.
    \item $\left\{\frac{n}{3n+ 1} : n \in \mathbb{N}\right\}$.
    \item $\left\{\frac{m}{m+ n} : m, n \in \mathbb{N}\right\}$.
\end{enumerate}
\end{exercise}

\sol{To avoid writing out every set definition, I am going to denote each set as $A_n$ where $n$ corresponds to the numerical value of the list from (a) - (d).
\begin{enumerate}
    \item $\sup A_1 = 1$, $\inf A_1 = 0$
    \item $\sup A_2 = 1$, $\inf A_2 = -1$
    \item $\sup A_3 = \frac{1}{3}$, $\inf A_3 = \frac{1}{4}$
    \item $\sup A_4 = 1$, $\inf A_3 = 0$
\end{enumerate}
}

\section{Consequences of Completeness}

\begin{ntheorem}
    {Nested Interval Property}For each $n \in \N$, assume we are given a closed interval $I_n = (a_n, b_n)$. Assume $I_n$ contains $I_{n+1}$. This results in a nested sequence of intervals. $$I_1 \supseteq I_2 \supseteq I_3 \supseteq I_4\dots$$ Then, $\bigcap^\infty_{n=1} \ne \emptyset$.
\end{ntheorem}

\textbf{tl;dr} there has to be something that is common to all of the sets.

\tpf{
    Notice that the sequence, $a_1,a_2,a_3,\dots$ is increasing. In other words, for each $n \in \N$, since $I_n \supset I_{n + 1}$ we have $a_n \leq a_{n+1}$. If we consider the set $A = \{a_n \ \colon \ n\in \N\}$. The element $b_1$ is an upper bound of $A$. (Note that $b_1$ and $a_1$ corresponds to the endpoints of the first set, $I_1$. Think of this as a tornado looking structure where the larger the $I_n$, the smaller the number line.) For each $n \in \N$, $a_n \leq b_n \leq b_1$. \\

    Since $A$ has an upper bound, it must have a least upper bound. Hence, let $\alpha = \sup A$. We claim that $\alpha \in \bigcap^\infty_{n=1} I_n$. We said $b_1$ was an upper bound. In fact, every $b_n$ is an upper bound of $A$. Choose any $n,m \in \N$. We want to show that $a_n \leq b_m$. Consider the following cases: \\

    \textbf{Case 1:} If $n < m$, then $a_n \leq a_m \leq b_m$. (Think: two number lines stacked on top of each other. The top number line is larger, call it $I_n$ and it has $a_n$ and $b_n$ as endpoints. Consider a contained line ($I_n \supseteq I_m$) that is smaller, and has endpoints $a_m$ and $b_m$.) \\

    \textbf{Case 2:} If $n > m$, then $a_n \leq b_n \leq b_m$. So every $b_n$ is an upper bound of $A$.\\

    Hence, 
    \begin{itemize}
        \item Because $\alpha = \sup A$, we have $\alpha \geq a_n$. 
        \item Since $b_n$ is an upper bound of $A$, we have $\alpha \leq b_n$.
    \end{itemize}
    so, $\alpha \in [a_n, b_n] = I_n$. Thus, $\alpha \in \bigcap^\infty_{n=1} I_n.$
}

Nested, closed, Bounded Intervals $\Rightarrow$ non-empty intersection.

\begin{ntheorem}
    {Archimedean Principle} 
    \begin{enumerate}
        \item Given any number $x \in R$, there exists an $n \in N$ satisfying $n > x$.
        \item Given any real number $y > 0$, there exists an $n \in N$ satisfying $1/n < y$.
    \end{enumerate}
\end{ntheorem}

\tpf{
    \begin{enumerate}
        \item If $\N$ was bounded, then we can let $s \in \N = \sup \N$. However, we know that there is always a higher number (e.g., $n+1$) for any $n \in \N$ that is given. Thus, by contradiction, there must exist $n \geq x$.
        \item For any $x > 0$, there exists $n \in \N$ such that $\frac{1}{n} < x$. 
    \end{enumerate}
}

\begin{theorem}
    Density of $\Q$ in $\R$. For any $a,b \in \R$ with $a < b$, there exists $q \in \Q$ such that $a < q < b$.
\end{theorem}

\tpf{
Since $ b - a > 0$, there exists $n \in \N$ such that $\frac{1}{n} < b - a$. From the \namrefthm{Archimedean Principle}, since $a\times n \in \R$, there exists $m \in \N$ such that $a\times n < m$. Let $m$ be there smallest such natural numbers (by the well ordered principle). Since $m$ is the smallest such natural number, it follows that $m - 1 \leq a\times n < m$. We then see that $a < \frac{m}{n}$. Now, we need to find some $\frac{m}{n} < b$. \begin{align*}
    m - 1 &\leq a\times n \\
    m &\leq a \times n + 1 \\
    \frac{m}{n} &\leq a + \frac{1}{n} \\
    \frac{m}{n} &< a + (b - a) \\
    \frac{m}{n} &< b
\end{align*} We now have that $a < \frac{m}{n} < b$ so $ \frac{m}{n}$ is a rational number in $(a,b)$
}

\begin{exercise}
    {1.4.1}Recall that $\mathbb{I}$ stands for the set of irrational numbers. 
    \begin{enumerate}
        \item Show that if $a,b \in \Q$, then $ab$ and $a + b$ are elements of $\Q$ as well.
        \item Show that if $a \in \Q$ and $t \in \mathbb{I}$, then $a + t \in \mathbb{I}$ and $at \in \mathbb{I}$ as long as $a \ne 0$.
        \item Part (a) can be summarized by saying that $\Q$ is closed under addition and multiplication. Is $\mathbb{I}$ closed under addition and multiplication? Given two irrational numbers $s$ and $t$, what can we say about $s + t$ and $st$? In other words, are there two irrational numbers that can be added and multiplied such that you get a number $x$ such that $x \notin \mathbb{I}$.
    \end{enumerate}
\end{exercise}

\sol{
\begin{enumerate}
    \item Let $a,b \in \Q$. This means there exists some $p,q,a,b \in \Z$ such that $$a = \frac{p}{q}$$ and $$b = \frac{a}{b}$$ where $q,b \ne 0$. The product of these numbers is $$ab = \frac{p}{q} \cdot \frac{a}{b} = \frac{pa}{qb}.$$ Since $pa,qb \in \Z$, $ab \in \Q$. The sum of these numbers is $$a + b = \frac{p}{q} + \frac{a}{b} = \frac{pb + aq}{qb}.$$ Since $pb + aq, qb \in \Z$, $a + b \in \Q$. 
    \item Let $a \in \Q$ and $t \in \I$. Assume, for contradiction, that $a + t \in \Q$. This would imply $t = (a + t) - a$ (because we can subtract $t + a$ from the original equation and rearrange terms). Since $a + t, a \in \Q$ their sum would be rational because the rational numbers are closed under addition. However, that would contradict the assumption that $t \in \I$. Hence, $a + t \in \I$.   
    \item For $\I$, it is not closed under addition and multiplication. Consider the following counterexample: $\sqrt{2} + (-\sqrt{2}) = 0$ which is not in the irrationals. For multiplication, consider $\sqrt{2} \cdot \sqrt{2} = 2$, which is also not in the irrationals. 
\end{enumerate}
}

\section{Cardinality}

Two sets have the same \textit{cardinality} if there exists a bijection between them. Thus, the natural numbers, the integers, and the rational numbers have the same cardinality. A set is \textit{countably infinite} if it has the same cardinality as $\N$. (If it can be put into one-to-one correspondence with $\N$.) A set is \textit{countable} if it is countably infinite or finite. \\

\setcounter{BoxCounter}{5}
\begin{theorem}
    $\R$ is not countable.
\end{theorem}

\tpf{%
    \textit{1} (most common) \\
    Suppose $\R$ is countable. Then we can list them all, or we can enumerate them. $\R = \{x_1,x_2,x_3,x_4,\dots\}$. We can write the decimal expansion of each of these. Consider the following table:
\begin{center}
\begin{tabular}{lllllll}
$x_1 =$ & $\boxed{a_{10}}$ & $a_{11}$ & $a_{12}$ & $a_{13}$ & $a_{14} $ & \dots \\ \hline
$x_2 =$ & $a_{20} $ & $\boxed{a_{21}}$ & $a_{22} $ & $a_{23} $ & $a_{24} $ & \dots \\ \hline
$x_3 =$ & $a_{30} $ & $a_{31} $ & $\boxed{a_{32}}$ & $a_{33} $ & $a_{34} $ & \dots \\ \hline
$x_4 =$ & $a_{40} $ & $a_{41} $ & $a_{42} $ & $\boxed{a_{43}}$ & $a_{44} $ & \dots \\ \hline
$x_5 =$ & $a_{50} $ & $a_{51} $ & $a_{52} $ & $a_{53} $ & $\boxed{a_{54}}$ & \dots \\ \hline
$x_6 =$ & $a_{60} $ & $a_{61} $ & $a_{62} $ & $a_{63} $ & $a_{64} $ & \dots
\end{tabular}
\end{center}

We will now construct a number that is not in this list. Focus on diagonal entries. For each $n \in \N$, let $b_n$ be a digit that is different fron $a_{nn}$. Now consider the number $y = 0.b_1b_2b_3b_4b_5\dots$. This number $y$ is not in our list. So our list did not include all of $\R$. Avoid repeating $9$s.
}

\tpf{
\textit{2} (uses nested interval theorem) \\
Suppose $\R$ is countable. Then we can enumerate $\R$ $\R = \{x_1,x_2,x_3,\dots\}$. Let $I_1$ be any closed interval that does not contain $x_1$. Next, we will find another closed interval $I_2$ that:\begin{itemize}
    \item $I_2 \subseteq I_1$ 
    \item $x_2 \notin I_2$
\end{itemize} Continue in this fashion creating a sequence of nested closed intervals: $I_1\supseteq I_2 \supseteq I_3 \supseteq \dots$ such that for all $k \in \N$, $x_k \notin I_k$. Now consider: $$\bigcap^\infty_{n=1} I_n$$
\begin{itemize}
    \item For each $k \in \N$, since $x_k \notin I_k$, we see $x_k \notin \bigcap^\infty_{n=1} I_n$.
    \item By the nested interval theorem, there exists $x \in \R$ such that $x \in \bigcap^\infty_{n=1} I_n$. So $x$ is a real number that is not included in our list.
\end{itemize}
}

\begin{theorem}
    A countable collection of finite sets is \textit{countable}.
\end{theorem}

\begin{theorem}
    \begin{enumerate}[label=(\roman*)]
        \item The union of two countable sets is \textit{countable}.
        \item A countable union of countable sets is \textit{countable}.
    \end{enumerate}
\end{theorem}

From Theorem \ref{thm:1.5.6}, we know that $\R$ is uncountable, but what about $(0,1)$? It does have the same cardinality of $\R$ because we can make a one-to-one and onto function between both the sets. Similarly, $(a,b)$ also has the same cardinality. What about $[a,b]$? 

\textbf{Recap:} $\N$ is countable, and $\R$ is uncountable and has a different cardinality than $\N$. Thus, the question is, do all uncountable sets have the same cardinality as $\R$? The answer is \textbf{no}.
    
\begin{ntheorem}
    {Canter's Theorem} For any set $A$, there does not exist an onto map from $A$ into $\mathcal{P}$.
\end{ntheorem}

\tpf{
    Suppose there exists an onto function, $f \ \colon \ A \rightarrow \mathcal{P}(A)$. So each $a \in A$ is mapped to an element $f(a) \in \mathcal{P}(A)$. Then, $f(a) \subseteq A$. We are going to construct an element of $\mathcal{P}(A)$ which is not mapped to by $f$. 

    Consider $B = \{a \in A \ \colon \ a \notin f(a)\}$. Since $f$ is onto there exists $a' \in A$ such that $B = f(a')$. Thus, there are two cases to consider:

    \begin{itemize}
        \item \textbf{Case 1:} If $a' \in B = f(a')$, then $a' \notin B$.
        \item \textbf{Case 2:} If $a' \notin B = f(a')$, then $a' \in B$.
    \end{itemize}

    As evidenced, both cases lead to contradictions, so $B$ is not the image of any $a \in A$. Therefore $f$ is not onto.
}

\begin{example}
    {Set and Power Set Matching}$A = \{a,b,c\}$. 
\end{example}

\sol{
$\mathcal{P}(A) = \emptyset, \{a\}, \{b\}, \{c\}, \{a,b\}, \{a,c\}, \{b,c\}, \{a,b,c\}$. Note that you can map $\{a\}$, $\{b\}$, $\{c\}$, to elements such as $\emptyset$, $\{a,b\}$, $\{a,b,c\}$, but there are still more elements that are left unmapped. We can extrapolate from our proof a set $B$ such that $B = \{a,c\}$ because those elements are not mapped to.
}

All of this is to show $\mathcal{P}(\R)$ has a larger cardinality than $\R$. Then $\mathcal{P}(\mathcal{P}(\R))$ has a larger cardinality than $\mathcal{P}(\R)$.

\stepcounter{chapter}{1}

\chapter{Sequences and Series}\label{sequences and series}
\vspace*{-0.25in}
\section{Discussion: Rearrangement of Infinite Series}

\noindent \textbf{Questions:}

\textbf{What is a \textit{sequence}?} \\

A countable, ordered list of elements. An example could be $1,2,3,4,5,\dots$. Note that this is \textit{ordered}, therefore distinguishing it from a sequence like $3,1,2,4,5,6,\dots$. Hence, order matters.

% This needs to be a part of the definition environment
A \textit{sequence} is a function whose domain is $\N$. \textbf{Note:} The domain $\N$ refers to each element's position in the list. For example, $(a_n) = a_1,a_2,a_3,\dots$.

We will focus on the \textit{limit} of a sequence. We use sequences to approximate other things.

\begin{example}
    {Sequence}$3,3.1,3.14,3.141,3.1415,\dots \approx \pi$. \\
    $x, x - \frac{x^3}{6}, x - \frac{x^3}{6} + \frac{x^5}{120}, x - \frac{x^3}{6} + \frac{x^5}{120} - \frac{x^7}{7!}, \dots \approx \sin(x)$
\end{example}

\noindent\textbf{What is a \textit{series}?} \\

An infinite sum. We look at the sequence of partial sums. We ask, do the partial sums approach a limit?

\begin{example}
    {Alternating Harmonic Series}$1 - 1/2 + 1/3 - 1/4 + 1/5 - 1/6 + \dots$
\end{example}

We can rearrange these terms such that we can `force' the series to converge to a specific number. Therefore, we will need to be careful with our definitions.



\section{The Limit of a Sequence}

\begin{definition}
    A \textit{sequence} is a function whose domain is $\N$. We write $(a_n) = a_1,a_2,a_3,\dots$.
\end{definition}
\setcounter{BoxCounter}{2}
\begin{definition}
    The sequence $(a_n)$ converges to $L$ if for all $\epsilon > 0$, there exists $N \in \N$ such that for all $n \geq \N$, $|a_n - L| < \epsilon$. In other words, there exists $N \in \N$ such that
    \begin{itemize}
        \item \textbf{(In the interval)} $a_N \in (L - \epsilon, L + \epsilon)$.
        \item \textbf{(Stays in the interval)} $\forall n \geq N$, $a_n \in (L - \epsilon, L + \epsilon)$.
    \end{itemize}
\end{definition}

\begin{example}
    {In-class}Let $a_n = \frac{1}{n}$. $(a_n) = (1,\frac{1}{2},\frac{1}{3},\dots$)
\end{example}

\epf{
    Our claim is $\displaystyle \lim_{n \rightarrow \infty} \frac{1}{n} = 0$. Thus, let $\epsilon = .01$. Does the sequence eventually get inside $(-.01,.01)$? We will set $N = 101$. So, for any $n \geq |0|$, $$\left|\frac{1}{n} - 0\right| = \frac{1}{n} \leq \frac{1}{101} < .01.$$ From $A_n$ and on, the sequence stayed within $\epsilon$ of $0$. But what about $\epsilon = .001$, $\epsilon = .00001$ and so on? \\

    Actual proof let $\epsilon > 0$. By the Archimedean Principle, there exists $N \in \N$ such that $N > \frac{1}{\epsilon}$. Now, for any $n \geq N$, $$\left|\frac{1}{n}\right| = \frac{1}{n} \leq \frac{1}{N} < \frac{1}{1/\epsilon}.$$ (Where $\frac{1}{1/\epsilon} = \epsilon$, but is in that form for demonstration purposes.) Therefore $\displaystyle \lim_{n \rightarrow \infty} \frac{1}{n} = 0$
}

``To get close'' means is that we are finding a bigger and bigger $N$ as $\epsilon$ gets smaller. Note that the choice of $N$ certainly depends on $\epsilon$.

\subsection{Basic Structure of a Limit Proof}

Claim: $\displaystyle \lim_{n \rightarrow \infty} a_n = L$. \\
Proof: Let $\epsilon > 0$. There exists $N \in \N$ such that \{something involving $\epsilon$\}. Assume $n \geq N$. Then, $$|a-n - L| \ \boxed{\dots} < \epsilon$$ (Where $\boxed{\dots}$ is going to be where the majority of the work is going to lie.

\newpage

\begin{example}
    {In-class}Claim: $\displaystyle \lim_{n \rightarrow \infty} \frac{2n-3}{2n} = 1$
\end{example}

\epf{
    Let $\epsilon > 0$. \textit{Scratch paper:} Solve for: $$\left|\frac{2n - 3}{2n} - 1\right| = \left|\frac{-3}{2n}\right| = \frac{3}{2n} < \epsilon \Rightarrow \frac{3}{2\epsilon} < n.$$ By the \namrefthm{Archimedean Principle}, there exists $N \in \N$ such that $N > \frac{3}{2\epsilon}$. Assume $n \geq N$, (want to know what happens past this point) $$\left|\frac{2n-3}{2n}-1\right| \leq \frac{3}{2N} < \frac{3}{2 \cdot 3/2\epsilon} = \epsilon.$$ Therefore, $\displaystyle \lim_{n \rightarrow \infty} \frac{2n-3}{2n} = 1$
}

\begin{example}
    Claim: $\limn{\frac{2n^2 + 1}{n^2}} = 2$
\end{example}

\epf{
    Let $\epsilon > 0$. By the \namrefthm{Archimedean Principle}, there exists $N \in \N$ such that [leave off] \textit{Scratch paper:} Solve for $$\disabs{\frac{2n^2 + 1}{n^2} - 2} =  \frac{2n^2}{n^2} < \epsilon \Rightarrow \frac{3}{2\epsilon} < n$$ [pick up] there exists $N \in \N$ such that $$N > \frac{1}{\sqrt{\epsilon}}.$$ Assume $n \geq N$, then
    \begin{align*}
        \disabs{\frac{2n^2 + 1}{n^2} - 2} & = \frac{1}{n^2}                    \\
                                          & \leq \frac{1}{N^2}                 \\
                                          & < \frac{1}{(1/(\sqrt{\epsilon})^2} \\
                                          & = \frac{1}{1/\epsilon}             \\
                                          & = \epsilon
    \end{align*}
    Therefore, $\limn{\frac{2n^2 + 1}{n^2}} = 2$
}

\begin{example}
    {In-class}Claim: $\limn{\frac{7n + 8}{3n + 6}} = \frac{7}{3}$
\end{example}

\epf{
    \begin{align*}
        \disabs{\frac{7n+8}{3n+6} - \frac{7}{3}} & = \disabs{\frac{21n + 24}{3(3n + 6)} - \frac{21n + 42}{3(3n + 6}} \\
                                                 & = \disabs{\frac{-18}{9n + 18}}                                    \\
                                                 & = \frac{18}{9n + 18} < \epsilon**                                 \\
                                                 & = \frac{18}{3} < 9n + 18                                          \\
                                                 & = \frac{18}{3} - 18 < 9n                                          \\
                                                 & = \frac{18/\epsilon - 18}{9} < n
    \end{align*}
    $** \frac{18}{9n + 8} < \frac{18}{9n} < \epsilon \Rightarrow \frac{2}{\epsilon} < N$. $\exists N \in \N$ such that $N > \frac{2}{\epsilon}$. Assume $n \geq N$, \begin{align*}
        \disabs{\frac{7n + 8}{3n + 6} - \frac{7}{3}} & = \frac{18}{9n + 18}   \\
                                                     & = \frac{2}{n+2}        \\
                                                     & < \frac{2}{n}          \\
                                                     & \leq \frac{2}{N}       \\
                                                     & < \frac{2}{\epsilon/2} \\
                                                     & = \epsilon
    \end{align*}
}

\textbf{Does every sequence have a limit?}

\begin{ntheorem}
    {Uniqueness of Limits}The limit when it exists, is unique.
\end{ntheorem}


\tpf{
    Let $(x_n)$ be a convergent sequence. Suppose $L$ and $M$ are limits of this sequence. Without the loss of generality, we are going to assume $M > L$ Let $$\epsilon = \frac{M - L}{3}.$$ Since $n_x$ converges to $L$, there exists $N_1 \in \N$ such that for all $n \geq N$, $\disabs{(x_n) - L} < \epsilon$. Since $(x_n)$ converges to $M$, there exists an $N_2 \in \N$ such that for all $n \geq N_2$, $\disabs{(x_n) - M} < \epsilon$. Consider $n = \max\{N_1, N_2\}$. Since $n \geq N_1$, $\disabs{(x_n) - L} < \epsilon$. Since $n \geq N_2$, $\disabs{(x_n) - M} < \epsilon$. Then $L - \epsilon < x_n < L + \epsilon$ and $M - \epsilon < x_n < M + \epsilon$. By our choice of $\epsilon$, we now have $$(x_n) < L + \epsilon < M - \epsilon < (x_n).$$ This is a contradiction. Thus, $(x_n)$ cannot have two different limits.
}

\begin{example}
    {}Let $(x_n) = \displaystyle \frac{\cos(n)}{3n}$. Claim: $\limn{(x_n)} = 0$
\end{example}

\epf{
    Let $\epsilon > 0$. By the \namrefthm{Archimedean Principle}, there exists $N \in \N$ such that $N > \frac{1}{3\epsilon}$ for all $n \geq N$,
    \begin{align*}
        \disabs{\frac{\cos(n)}{3n} - 0} & = \disabs{\frac{\cos(n)}{3n}} \\
                                        & \leq \frac{1}{3n}             \\
                                        & \leq \frac{1}{3N}             \\
                                        & < \frac{1}{3(1/ 3\epsilon)}   \\
                                        & = \epsilon
    \end{align*}
}

\begin{example}
    {}Let $(y_n) = \frac{4n-1}{n^2}$. Claim: $\limn{y_n} = 0$.
\end{example}

\epf{
    Let $\epsilon > 0$. By the \namrefthm{Archimedean Principle}, there exists $N \in \N$ such that $N > \frac{1}{\epsilon}$. For all $n \geq N$,
    \begin{align*}
        \disabs{\frac{4n-1}{n^2} - 0} & = \disabs{\frac{4n-1}{n^2}} \\
                                      & = \frac{4n-1}{n}            \\
                                      & < \frac{4n}{n^2}            \\
                                      & = \frac{4}{n}               \\
                                      & \leq \frac{4}{N}            \\
                                      & < \frac{4}{4 / \epsilon}    \\
                                      & = \epsilon
    \end{align*}
}

\section*{Exercises}

\begin{exercise}
    {2.2.2(b)}Verify, using \defref{2.2.3}, that
    the following sequences converge to the proposed limit.
    \begin{enumerate}
        \setcounter{enumi}{1}
        \item \(\limn{\frac{2n^2}{n^3+3}} = 0\)
    \end{enumerate}
\end{exercise}

\expf{\hfill
    \begin{enumerate}
        \setcounter{enumi}{1}
        \item Let \(\epsilon > 0\). By the \namrefthm{Archimedean Principle}, there exists an \(N \in \N\) such that \(N > \frac{2}{\epsilon}\). Then, for \(n \geq N\),
              \begin{align*}
                  \disabs{\frac{2n^2}{n^3 + 3} - 0} & = \disabs{\frac{2n^2}{n^3 + 3}} \\
                                                    & = \frac{2n^2}{n^3 + 3}          \\
                                                    & < \frac{2n^2}{n^3}              \\
                                                    & = \frac{2}{n}                   \\
                                                    & \leq \frac{2}{N}                \\
                                                    & = \frac{2}{2/\epsilon}          \\
                                                    & = \epsilon.
              \end{align*}
              Therefore, \(\limn{\frac{2n^2}{n^3 + 3}} = 0\). \qedhere
    \end{enumerate}
}

\begin{exercise}
    {2.2.3} Describe what we would have to demonstrate in order to disprove each of the following statements.
    \begin{enumerate}
        \item At every college in the United States, there is a student who is at least
              seven feet tall.
        \item For all colleges in the United States, there exists a professor who gives every student a grade of either A or B.
        \item There exists a college in the United States where every student is at least six feet tall.
    \end{enumerate}
\end{exercise}

\sol{
    \begin{enumerate}
        \item There is at least one college in the United States where all students are less than seven feet tall.
        \item There is at least one college in the United States where all professors give at least one student a grade of C or lower.
        \item For all colleges in the United States, there exists a student who is less than six feet tall.
    \end{enumerate}
}

\begin{exercise}
    {2.2.4} Give an example of each or state that the request is impossible.
    For any that are impossible, give a compelling argument for why that is the case.
    \begin{enumerate}
        \item A sequence with an infinite number of ones that does not converge to one.
        \item A sequence with an infinite number of ones that converges to a limit not equal to one.
        \item A divergent sequence such such that for every \(n \in \N\) it is possible to find \(n\) consecutive ones somewhere in the sequence.
    \end{enumerate}
\end{exercise}

\sol{
    \begin{enumerate}
        \item Possible. Consider the piecewise function: \(a_n = \begin{cases}
                    1 & \text{if } n \text{ is even} \\
                    0 & \text{if } n \text{ is odd}
        \end{cases}\)
        \item Impossible. A sequence that converges must have its terms approach a specific value (the limit). If the sequence has an infinite number of ones, it must have subsequences of ones arbitrarily far out. For the sequence to converge to a limit different from 1, the terms would have to approach that different limit, say \(L \ne 1\), meaning the ones must become rare or eventually stop appearing, contradicting the infinite number of ones. Therefore, such a sequence is impossible.
        \item Possible. \((0,1,0,1,1,0,1,1,1,0,1,1,1,1,0,\dots)\)
    \end{enumerate}
}

\section{The Algebraic and Order Limit Theorems}

\begin{definition}
    A sequence $(x_n)$ is \textit{bounded} if there exists some $M > 0$ such that every term in the sequence belongs to $[-M,M]$.
\end{definition}

\begin{theorem}
    Every convergent sequence is bounded.
\end{theorem}

\tpf{
Let $(x_n)$ be a convergent sequence with limit $L$. There exists $N \in \N$ such that for all $n \geq N$, $\disabs{(x_n) - L} < 1$. Equivalently, $(x_n) \in (L - 1, L + 1)$. Let $$M = \max\{|x_1|,|x_2|,\dots,|x_{N-1}|,|L + 1|, |L - 1|\}.$$ We claim that for all $n \in \N$, $|x_n| \leq M$.
\begin{enumerate}
    \item This is true for $n < N$.
    \item For $n \geq N$, we know $L - 1 < x_n < L + 1$, so $(x_n) \leq \max\{| L - 1 |, |L + 1|\}$
\end{enumerate}
Thus, every term is in $[-M,M]$.
}

\begin{ntheorem}
    {Algebraic Limit Theorem}Let \(\limn{a_n} = a\) and \(\limn{b_n} = b\). Then,
    \begin{enumerate}[label=(\roman*)]
        \item \(\limn{ca_n} = ca\) for all $c \in \R$;
        \item \(\limn{(a_n + b_n)} = a + b\);
        \item \(\limn{(a_nb_n)} = ab\);
        \item \(\limn{\frac{a_n}{b_n}} = \frac{a}{b}\) provided $b \neq 0$.
    \end{enumerate}
\end{ntheorem}

\textit{Scratch Paper:}
\begin{align*}
    \disabs{ca_n - ca} = \disabs{c}\disabs{a_n - a} & < \epsilon                    \\
    \disabs{a_n - a}                                & < \frac{\epsilon}{\disabs{c}}
\end{align*}
\text{Leave off and go back to proof$^1$}
\tpf{(i) \\
Let $\epsilon > 0$.$^1$ Since $(a_n)$ converges to $a$, there exists $N \in \N$ such that for all $n \geq N$, $\disabs{a_n - a} < \frac{\epsilon}{\disabs{c}}$. Now, for any $n \geq N$ we have two case because we want to avoid dividing by 0:
\begin{itemize}
    \item If $c = 0$: \\
          then each $ca_n = 0$. So $(ca_n)$ converges to 0, which can equal $ca$.
    \item If $c > 0$: \\
          $\disabs{ca_n - ca} = \disabs{c}\disabs{a_n - a} < \disabs{c}\frac{\epsilon}{\disabs{c}} = \epsilon$. \qedhere
\end{itemize}

(ii) \\
\textit{Scratch paper:}
\begin{align}
    \disabs{(a_n + b_n)} & = \disabs{(a_n - a) + (b_n - b)}          \\
                         & \leq \disabs{a_n - a} + \disabs{b_n - b}  \\
                         & < \frac{\epsilon}{2} + \frac{\epsilon}{2}
\end{align}
Note that (2.2) is from the triangle inequality. Now, we will pick up to back at $\epsilon > 0$. \\

Let \(\epsilon > 0\). Since $(a_n)$ converges to $a$, there exists $N_1 \in \N$ such that for all $n \geq N_1$, $\disabs{a_n - a} < \frac{\epsilon}{2}$. Since $(b_n)$ converges to $b$, there exists $N_2 \in \N$ such that for all $n \geq N_2$, $\disabs{b_n - b} < \frac{\epsilon}{2}$. Now, let $N = \max\{N_1, N_2\}$. Thus, for any $n \geq N$, (refer back to scratch paper). \\

(iii) \\
\textit{Scratch paper:}
\begin{align}
    \disabs{a_nb_n - ab} & = \disabs{a_nb_n - ab_n + ab_n - ab}                                    \\
                         & = \disabs{a_n(b_n - b) + b(b_n - b)}                                    \\
                         & \leq \disabs{a_n}\disabs{b_n - b} + \disabs{b}\disabs{b_n - b}          \\
                         & \leq M\disabs{b_n - b} + M\disabs{a_n - a}.                             \\
                         & < M\left(\frac{\epsilon}{2M}\right) + M\left(\frac{\epsilon}{2M}\right) \\
                         & = \epsilon
\end{align}
Note that: (2.4) is where we added 0, (2.5) is from the triangle inequality, and (2.6) is just factored. Additionally, we choose \(N\) to get the fractions in (2.8) Now, we will pick up to back at $\epsilon > 0$. \\

Let \(\epsilon > 0\). Since convergent sequences are bounded, then there exists $M > 0$ such that for all $n \in \N$, $|a_n| \leq M$. We can choose \(M\) so that \(|b_n| \leq M\) as well. Since $(a_n)$ converges to $a$, there exists $N_1 \in \N$ such that for all $n \geq N_1$, $\disabs{a_n - a} < \frac{\epsilon}{2M}$. Since $(b_n)$ converges to $b$, there exists $N_2 \in \N$ such that for all $n \geq N_2$, $\disabs{b_n - b} < \frac{\epsilon}{2M}$. Now, let $N = \max\{N_1, N_2\}$. Thus, for any $n \geq N$, (refer back to scratch paper, and change (2.4)'s sign from an `=' to `$\leq$'). \\

(iv) \\
\textit{Scratch paper:}
\begin{align*}
    \disabs{\frac{a_n}{b_n} - \frac{a}{b}} & = \disabs{\frac{a_nb - ab_n}{b_nb}}                                                     \\
                                           & = \disabs{\frac{a_nb - ab_n + ab_n - ab}{b_nb}}                                         \\
                                           & = \disabs{\frac{a_n(b - b_n) + b(b_n - b)}{b_nb}}                                       \\
                                           & = \disabs{\frac{a_n(b - b_n)}{b_nb} + \frac{b(b_n - b)}{b_nb}}                          \\
                                           & \leq \disabs{\frac{a_n}{b_n}}\disabs{b - b_n} + \disabs{b}\disabs{\frac{b_n - b}{b_nb}} \\
                                           & < \epsilon
\end{align*}
Let \(\epsilon > 0\). Since \((b_n)\) converges to \(b\), there exists \(N_1 \in \N\) such that for all \(n \geq N_1\), \(\disabs{b_n } > \disabs{\frac{b}{2}}\). There also exists \(N_2 \in \N\) such that for all \(n \geq N_2\), \(\disabs{b_n - b} < \frac{\epsilon|b|^{2}}{2}\). Now, let \(N = \max\{N_1, N_2\}\). Let \(n \geq N\), (refer back to scratch paper). \qedhere
}

\begin{lemma}
    Let \((a_n)\) and \(c < a\). There exists \(N \in \N\) such that for all \(n \geq N\), \(a_n > c\). Similarly, if \(a < d\), there exists \(N \in \N\) such that for all \(n \geq N\), \(a_n < d\).
\end{lemma}

\subsection{Limits and Order}

\begin{ntheorem}
    {Order Limit Theorem}Let \((a_n)\) and \((b_n)\) be sequences. If \(\limn{a_n} = a\) and \(\limn{b_n} = b\), then
    \begin{enumerate}[label=(\roman*)]
        \item If \(a_n \geq c\) for all \(n \in \N\), then \(a \geq c\).
        \item If \(a_n \leq b_n\) for all \(n \in \N\), then \(a \leq b\).
        \item If there exists \( c \in \mathbb{R} \) for which \( c \leq b_n \) for all \( n \in \mathbb{N} \), then \( c \leq b \). Similarly, if \( a_n \leq c \) for all \( n \in \mathbb{N} \), then \( a \leq c \).
        
    \end{enumerate}
\end{ntheorem}



\section*{Exercises}

\begin{exercise}
    {2.3.1}
    \begin{enumerate}
        \item Assume \(\limn{x_n = 0}\) with \(x_n \geq 0\). Show that \(\limn{\sqrt{x_n}} = 0\).
        \item Assume \(\limn{x_n = 49}\) with \(x_n \geq 0\). Show that \(\limn{\sqrt{x_n} = 7}\)
    \end{enumerate}
\end{exercise}

\expf{ \hfill
    \begin{enumerate}
        \item Let \(\epsilon > 0\). Since \((x_n)\) converges to 0, there exists \(N \in \N\) such that for all \(n \geq N\), \(\disabs{x_n - 0} < \epsilon^{2}\). Now, for any \(n \geq N\), \(\disabs{\sqrt{x_n} - 0} = \sqrt{x_n} < \sqrt{\epsilon^{2}} = \epsilon\). Therefore, \(\limn \sqrt{x_n} = 0\).
        \item Let \(\epsilon > 0\). Since \((x_n)\) converges to 49, there exists \(N \in \N\) such that for all \(n \geq N\),
              \begin{align*}
                  \disabs{x_n - 7} & = \disabs{\frac{(x_n - 7)(x_n + 7)}{\sqrt{x_n} + 7}} \\
                                   & = \disabs{\frac{x_n - 49}{\sqrt{x_n} + 7}}           \\
                                   & \leq \frac{\disabs{x_n - 49}}{7}
              \end{align*}
    \end{enumerate}
}

\begin{exercise}
    {2.3.2}Using only \defref{2.2.3}, prove that if $(x_n) \rightarrow 2$, then
    \begin{enumerate}
        \item $\displaystyle \left(\frac{2x_n - 1}{3}\right) \rightarrow 1$;
        \item $\displaystyle \left(1 / x_n\right) \rightarrow 1 / 2$.
    \end{enumerate}
    (For this exercise the Algebraic Limit Theorem is off-limits, so to speak.)
\end{exercise}

\sol{ \hfill
    \begin{enumerate}
        \item \textit{Proof.} Let $\epsilon > 0$. Since $(x_n)$ converges to 2, there exists $N \in \N$ such that for all $n \geq N$, $\disabs{x_n - 2} < \epsilon$. Now, for any $n \geq N$,
              \begin{align*}
                  \disabs{\frac{2x_n - 1}{3} - 1} & = \disabs{\frac{2x_n - 1 - 3}{3}} \\
                                                  & = \disabs{\frac{2x_n - 4}{3}}     \\
                                                  & = \frac{2}{3}\disabs{x_n - 2}     \\
                                                  & < \disabs{x_n - 2}                \\
                                                  & < \epsilon
              \end{align*}
              Therefore, \(\frac{2x_n - 1}{3} \rightarrow 1\) \qed
        \item \textit{Proof.} Let $\epsilon > 0$. Since \((x_n)\) converges to 2, there exists \(N_1 \in \N\) such that for all \(n \geq N_1\), \(x_n \geq 1\). Then, we will choose \(N_2\) so that \(|x_n - 2 | < \epsilon\) for all \(n \geq N_2\). Afterwards, we take \(N = \max\{N_1, N_2\}\). And note that for \(n \geq N\),

              \begin{align*}
                  \disabs{\frac{1}{x_n} - \frac{1}{2}} & = \disabs{\frac{2 - x_n}{2x_n}}  \\
                                                       & < \frac{\disabs{2 - x_n}}{2}     \\
                                                       & < \frac{\epsilon}{2}             \\
                                                       & < \epsilon
              \end{align*} \qed
    \end{enumerate}
}

\section{The Monotone Convergence Theorem and a First Look at Infinite Series}

\begin{definition}
    A sequence \(a_n\) is \textit{increasing} if \(a_n \leq a_{n+1}\) for all \(n \in \N\) and \textit{decreasing} if \(a_n \geq a_{n+1}\) for all \(n \in \N\). A sequence is \textit{monotone} if it is either increasing or decreasing.
\end{definition}

\begin{ntheorem}
    {Monotone Convergence Theorem}If a sequence is monotone and bounded, then it converges.
\end{ntheorem}

\tpf{
    Let \((a_n)\) be an increasing and bounded sequence. Since \((a_n)\) is bounded, the set \(A = \{a_n \ |\ n \in \N\}\) is clearly also bounded. Since \(A\) is bounded, \(\sup{A}\) exists. We claim that \(\limn{a_n} = \sup{A}\). Thus, for all \(\epsilon > 0\) and by our definition of supremum, there exists \(N \in \N\) such that \(\sup{A} - \epsilon < a_N \leq \sup{A}\). Since \((a_n)\) is increasing, for all \(n \geq N\), \(\sup{A} - \epsilon < a_N \leq a_n \leq \sup{A}\). It follows that \(\disabs{a_n - \sup{A}} < \epsilon\). Therefore, \(\limn{a_n} = \sup{A}\).
}

\begin{example}
    {MCT}Consider the recursively defined sequence \(x_n\) where \(x_1 = 3\) and for all \(n \in \N\), \(x_{n + 1} = \frac{1}{4 - x_n}\). Show that \(x_n\) converges.
\end{example}

\epf{
    We will show that \(x_n\) is monotone and bounded.
    \begin{itemize}
        \item \textbf{Part 1: Monotone Decreasing}
              \begin{itemize}
                  \item \underline{Base case}: \(x_1 = 3\), \(x_2 = 1\).
                  \item \underline{Induction step}: Assume for some \(n \in \N\), \(x_n \geq x_{n+1}\). It follows that
                        \begin{align*}
                            x_n               & \geq x_{n + 1}               \\
                            4 - x_n           & \leq 4 - x_{n + 1}           \\
                            \frac{1}{4 - x_n} & \geq \frac{1}{4 - x_{n + 1}} \\
                            x_{n + 1}         & \geq x_{n + 2}
                        \end{align*}
              \end{itemize}
        \item \textbf{Part 2: Bounded Below}
              Claim: Sequence is bounded below by 0.
              \begin{itemize}
                  \item \underline{Base case}: \(x_1 = 3 > 0\).
                  \item \underline{Induction step}: Assume for some \(n \in \N\), \(x_n \geq 0\). It follows that \(4 - x_n \leq 4\), and when we take the reciprocal, we get \begin{align*}
                            \frac{1}{4 - x_n} & \leq \frac{1}{4} \\
                            x_{n+1}           & \geq 1/4         \\
                                              & > 0
                        \end{align*}
                        By math induction, \(x_n\) is bounded below by 0.
              \end{itemize}
    \end{itemize}
    By the Monotone Convergence Theorem, \(x_n\) converges.

    So, what is the limit? We know \((x_n)\) converges so let \(L = \limn{x_n}\). Then, \(\limn{x_{n+1}} = L\). We also know \(x_{n+1} = \frac{1}{4 - x_n}\). So \(L = \limn{x_{n+1}} = \limn{\frac{1}{4- x_n} = \frac{1}{4 - L}}\). It must be true that \(L = \frac{1}{4 - L}\). Solving for \(L\), we get \begin{align*}
        L(4 - L)     & = 1 \\
        4L - L^2     & = 1 \\
        L^2 - 4L + 1 & = 0
    \end{align*} Hence, \(L = 2 - \sqrt{3}\) or \(L = 2 + \sqrt{3}\). Notice that it cannot be the latter because it is bigger than 3.
}

\subsection{Recap and Summary}

We use limits to define multiple things in calculus. This is why we are focusing so heavily upon it. For example,

\begin{enumerate}
    \item Derivatives: \(\lim_{h \to 0} \frac{f(x + h) - f(x)}{h}\)
    \item Integrals: \(\lim_{n \to \infty} \sum_{i = 1}^{n} f(x_i) \Delta x\)
    \item Infinite Series: \(\lim_{n \to \infty} \sum_{i = 1}^{n} a_i\)
\end{enumerate}

Looking to the future, we are going to use functions and summations together. For example, when we have \(f(x) = \sum_{n=0}^\infty (a_n)x^n\) such that \(f'(x) = \sum_{n=0}^\infty (a_n)x^{n-1}\).

% Formalize this later:

\textit{Limit Inferior:} Let \(x_n\) be a bounded sequence, then \(\liminf_{n \to \infty} x_n = \lim_{n \to \infty} \inf\{x_k \ | \ k \geq n\}\). This is the largest a limit can get.

\textit{Limit Superior:} Let \((x_n)\) be a bounded sequence, then \(\limsup_{n \to \infty} x_n = \lim_{n \to \infty} \sup\{x_k \ | \ k \geq n\}\). This is the smallest a limit can get.

\textbf{See Exercise 2.4.7 in the book for more information.}

\begin{example}
    {\textit{Monotone Decreasing Sequence}} \begin{align*}
        x_1, x_2, x_3, x_4, x_5, x_6, &\dots \sup\{x_k \mid k \geq 1\} = S. \\
        x_2, x_3, x_4, x_5, x_6, &\dots \sup\{x_k \mid k \geq 2\} = S. \\
        x_3, x_4, x_5, x_6, &\dots \sup\{x_k \mid k \geq 3\} = S. \\
        x_4, x_5, x_6, &\dots \sup\{x_k \mid k \geq 4\} = S. 
    \end{align*}
\end{example}

\(\limsup_{n \to \infty} x_n\) is guaranteed to exist by the \namrefthm{Monotone Convergence Theorem}. 

\begin{example}
    {liminf} Let \(x_n = (-1)^n(1 + \frac{1}{n})\). Thus, \(x_{1,2,3} = -2,1\frac{1}{2},-1\frac{1}{3}\dots\)
\end{example}

\begin{example}
    {Convergence Towards 0}Let \(x_n = (-1)^n\frac{1}{n}\). Thus, \(x_{1,2,3} = -1, \frac{1}{2}, -\frac{1}{3}\dots\).
\end{example}

\begin{theorem}
    A sequence \(x_n\) is convergent if, and only if, \(\liminf_{n \to \infty} x_n = \limsup_{n \to \infty} x_n\).  
\end{theorem} See Theorem 2.4.6 in the book for another view.

\section{Subsequences and the Bolzano-Weierstrass Theorem}

\begin{definition}
    Let \(a_n\) be a sequence of real numbers, and let \(n_1 < n_2 < n_3 < \dots\) be an increasing sequence of natural numbers. Then, the sequence \(a_{n_1}, a_{n_2}, a_{n_3}, \dots\) is called a \textit{subsequence} of \(a_n\) and is denoted by \(a_{n_k}\), where \(k \in \N\) indexes the subsequence.
\end{definition}

\begin{theorem}
    Subsequences of a convergent sequence converge to the same limit as the original sequence.
\end{theorem}

\tpf{
    Let \(x_{n_k}\) be a subsequence of \(x_n\), and let \(L = \limn{x_n}\). We want to show that \(\limn{x_{n_k}} = L\). Let \(\epsilon > 0\). Since \(x_n\) converges to \(L\), there exists \(N \in \N\) such that for all \(n \geq N\), \(\disabs{x_n - L} < \epsilon\). Since \(n_k\) is increasing, there exists \(M \in \N\) such that \(n_k \geq N\) for all \(k \geq M\). Thus, for all \(k \geq M\), \(\disabs{x_{n_k} - L} < \epsilon\). Therefore, \(\limn{x_{n_k}} = L\).


    Let \(x_{n_k}\) be a subsequence of \(x_n\). Let \(\epsilon > 0\). Since \((x_n) \rightarrow L\), there exists an \(N \in \N\) such that for all \(n \geq N\), \(\disabs{x_n - L} < \epsilon\).

    Now, looking at \(x_{n_k}\), notice that \(n_k \geq k\) for all \(k\). Consider \(k = N\). For any \(n \geq N\), \(n \geq N \geq k\). Thus, \(\disabs{x_{n_k} - L} < \epsilon\). Therefore, \(\limn{x_{n_k}} = L\).
}

\begin{ntheorem}
    {Divergence Criterion}If \(x_n\) has two subsequences that converge to different limits, then \(x_n\) diverges.
\end{ntheorem}

\begin{ntheorem}
    {Bolzano-Weierstrass Theorem}Every bounded sequence in \(\R\) has a convergent subsequence.
\end{ntheorem}

\tpf{
    Let \(x_n\) be a bounded sequence. There exists an \(M > 0\) such that every term \(x_n\) belongs to \([-M,M]\). To prove this theorem, we will be utilizing a recursive argument style. Thus, let \(I_0 = [-M,M]\). \(I_0\) has length \(2M\). Cut \(I_0\) in half with \(I_1\) and \(I_2\) both being half as long as \(I_0\). Since \(x_n\) is bounded, there exists an \(I_L\) or \(I_R\) that contains infinitely many terms of \(x_n\). We will pick one, call it \(I_1\) that is contained in \(I_L\). \(I_1\) has length \(M\). Pick one of those terms inside \(I_1\) and call it \(x_{n_1}\). Now, cut \(I_1\) in half with equal length in intervals. One of them contains infinitely many terms. Call that interval \(I_2\). \(I_2\) has length \(\frac{M}{2}\). Pick one of those terms inside \(I_2\) and call it \(x_{n_2}\). Continue this process indefinitely for all \(n \N\) with \(n_1 > n_2\). Continue this process, and we get 
    \begin{itemize}
        \item a sequence of closed intervals \(I_n\).
        \begin{itemize}
            \item \(I_n\) has length \(\frac{2M}{2^n}\).
            \item They are nested, \(I_n \subseteq I_{n - 1}\).
        \end{itemize}
            \item a subsequence \(x_{n_k}\)
            \begin{itemize}
                \item for all \(k_1\), \(x_{n_k} \in I_k\).
            \end{itemize}
    \end{itemize}



    The \namrefthm{Nested Interval Property} states that \(\displaystyle \bigcup_{n = 1}^\infty I_n\) is non empty. Let \(L\) be a point in \(\displaystyle \bigcup_{n = 1}^\infty I_n\). We claim \(\limn{x_{n_k}} = L\). Let \(\epsilon > 0\). There exists an \(N \in \N\) such that \(\frac{2M}{2^n} < \epsilon\). (Since \(\limn{\frac{2M}{2^n}} = 0\). See ) For any \(k \geq N\), recall that \(x_{n_k}\), \(L \in I_k\). Since \(I_k\) has length \(\frac{2M}{2^n}\). Thus, \(\disabs{x_{n_k} - L} < \epsilon\). Therefore, \(\limn{x_{n_k}} = L\) and \((x_n)\) has a convergence sequence. 
}

\begin{tikzpicture}

    % Define the interval length
    \def\M{6} % Length of I_0
    
    % Draw the first interval I_0
    \draw[thick] (-\M, 0) -- (\M, 0) node[anchor=west] {$I_0 = [-M,M]$};
    \draw[thick] (-\M, 0.1) -- (-\M, -0.1) node[below] {$-M$};
    \draw[thick] (\M, 0.1) -- (\M, -0.1) node[below] {$M$};
    
    % Draw the midpoint of I_0
    \draw[thick] (0, 0.1) -- (0, -0.1) node[below] {$0$};
    
    % Draw subintervals I_1, I_L, I_R
    \def\halfM{\M / 2}
    \draw[thick] (-\halfM, -1) -- (\halfM, -1) node[anchor=west] {$I_1$};
    \draw[thick] (-\halfM, -0.9) -- (-\halfM, -1.1) node[below] {};
    \draw[thick] (\halfM, -0.9) -- (\halfM, -1.1) node[below] {};
    
    % Left and Right halves
    \draw[dashed] (-\M, -1.5) -- (0, -1.5) node[anchor=west] {$I_L$};
    \draw[dashed] (0, -1.5) -- (\M, -1.5) node[anchor=west] {$I_R$};
    
    % Draw I_2 within I_1
    \def\quarterM{\M / 4}
    \draw[thick] (-\quarterM, -2) -- (\quarterM, -2) node[anchor=west] {$I_2$};
    \draw[thick] (-\quarterM, -1.9) -- (-\quarterM, -2.1) node[below] {};
    \draw[thick] (\quarterM, -1.9) -- (\quarterM, -2.1) node[below] {};
    
    % Further subintervals
    \def\eighthM{\M / 8}
    \draw[thick] (-\eighthM, -3) -- (\eighthM, -3) node[anchor=west] {$I_3$};
    \draw[thick] (-\eighthM, -2.9) -- (-\eighthM, -3.1) node[below] {};
    \draw[thick] (\eighthM, -2.9) -- (\eighthM, -3.1) node[below] {};
    
    % Points representing x_n's
    \filldraw[red] (1.5, 0) circle (2pt) node[anchor=north] {$x_{n_1}$};
    \filldraw[red] (0.7, -1) circle (2pt) node[anchor=north] {$x_{n_2}$};
    \filldraw[red] (0.35, -2) circle (2pt) node[anchor=north] {$x_{n_3}$};
    \filldraw[red] (0.18, -3) circle (2pt) node[anchor=north] {$x_{n_4}$};
    
\end{tikzpicture}

\begin{theorem}
    Let \(b \in (0,1)\). Then \(\limn{b^n = 0}\).
\end{theorem}

\tpf{
    The sequence \((b^n)\) is monotone decreasing. This is because \(b^{n + 1} = b^nb < b^n\). This sequence is also bounded by 0. Hence, by the \namrefthm{Monotone Convergence Theorem}, \((b^n)\) converges. Now, let \(L = \limn{b^n}\). Consider the subsequence \(b^{2n}\). This sequence also converges to \(L\). Thus, \begin{align}
        L & = \limn{b^{2n}} \\
          & = \limn{b^nb^n} \\
          & = \limn{b^n}\limn{b^n} \\
            & = L^2.
    \end{align} Thus, \(L = 0\) or \(L = 1\). The limit cannot be 1 because \(b^n\) is decreasing away from 1. Therefore, \(L = 0\).
}

% % \chapter{Equivalency} \label{Equivalency}
% % \vspace*{-0.25in}
% % \input{Chapters/ch_7}

% % \chapter{Functions \& Bijections}\label{Functions}
% % \vspace*{-0.25in}
% % \input{Chapters/ch_8}

% % \chapter{Recursion \& The Whole Numbers}\label{Recursion}
% % \vspace*{-0.25in}
% % \input{Chapters/ch_9}

% % \chapter{The Integers} \label{Integers}
% % \vspace*{-0.25in}
% % \input{Chapters/ch_10}

% % \chapter{Cardinality} \label{Cardinality}
% % \vspace*{-0.25in}
% % \input{Chapters/ch_11}

% \chapter*{Presentation Problem 1}\label{pres 1}
% \vspace*{0.25in}
% 
\begin{exercise}
    {1.2.13}For this exercise, assume \hyperref[ex:1.2.5]{Exercise 1.2.5} has been successfully completed.
    \begin{enumerate}[label=(\alph*)]
        \item Show how induction can be used to conclude that
        \[
        (A_1 \cup A_2 \cup \cdots \cup A_n)^c = A_1^c \cap A_2^c \cap \cdots \cap A_n^c
        \]
        for any finite \(n \in \mathbb{N}\).
    
        \item It is tempting to appeal to induction to conclude
        \[
        \left( \bigcup_{i=1}^{\infty} A_i \right)^c = \bigcap_{i=1}^{\infty} A_i^c,
        \]
        but induction does not apply here. Induction is used to prove that a particular statement holds for every value of \(n \in \mathbb{N}\), but this does not imply the validity of the infinite case. To illustrate this point, find an example of a collection of sets \(B_1, B_2, B_3, \dots\) where
        \[
        \bigcap_{i=1}^{n} B_i \neq \emptyset \quad \text{is true for every } n \in \mathbb{N},
        \]
        but
        \[
        \bigcap_{i=1}^{\infty} B_i = \emptyset
        \]
        fails.
    
        \item Nevertheless, the infinite version of De Morgan’s Law stated in (b) is a valid statement. Provide a proof that does not use induction.
    \end{enumerate}
\end{exercise}
\newpage
\sbseteqpf%
%
[In this proof, we plan to prove (c). Thus, we need to show that:
\[
\left( \bigcup_{i=1}^{\infty} A_i \right)^c \subseteq \bigcap_{i=1}^{\infty} A_i^c
\]
and
\[
\left( \bigcup_{i=1}^{\infty} A_i \right)^c \supseteq \bigcap_{i=1}^{\infty} A_i^c.
\]]%
%
{Let \(x \in \left( \bigcup_{i=1}^{\infty} A_i \right)^c\). This means \(x\) is in the union set of \(A_i\) for all \(i \in \N\). Then, because we are taking the complement of \(\left( \bigcup_{i=1}^{\infty} A_i \right)\), that means \(x \notin A_i\) for all \(i \in \mathbb{N}\). Hence, \(x\) is in the complement of each \(A_i\). Thus, we can use the definition of intersection to assert \(x \in \bigcap_{i=1}^{\infty} A_i^c\).

Therefore, we have shown:
\[
\left( \bigcup_{i=1}^{\infty} A_i \right)^c \subseteq \ \bigcap_{i=1}^{\infty} A_i^c.
\]}%
%
%
%
{Similar to before, let \(x \in \bigcap_{i=1}^{\infty} A_i^c\). Because \(x \in A^c_i\) for all \(i \in \N\) we know \(x \notin A_i\). Hence, \(x \notin \left(\bigcup_{i=1}^{\infty} A_i\right)\), which means \(x \in \left( \bigcup_{i=1}^{\infty} A_i \right)^c\).

Therefore, we have shown:
\[
\left( \bigcup_{i=1}^{\infty} A_i \right)^c \supseteq \ \bigcap_{i=1}^{\infty} A_i^c.
\]}%
%
%
{By showing both inclusions, we see that:
\[
\left( \bigcup_{i=1}^{\infty} A_i \right)^c = \bigcap_{i=1}^{\infty} A_i^c.
\] \qedhere
}%
%
%
{xgray}


% \addcontentsline{toc}{chapter}{Presentation Problem 1}

\printglossary[title=Glossary]

%-%-%-%-%-%-%-%-%-%-%-%-%-%-%-%-%-%-%-%-%-%-%-%-%-%-%-%-%-%-%-%-%-%-%-%-%-%-%

\end{document}
%-%-%-%-%-%-%-%-%-%-%-%-%-%-%-%-%-%-%-%-%-%-%-%-%-%-%-%-%-%-%-%-%-%-%-%-%-%-%

%-%-%-%-%-%-%-%-%-%-%-%-%-%-%-%-%-%-%-%-%-%-%-%-%-%-%-%-%-%-%-%-%-%-%-%-%-%-%